\documentclass[11pt,a4paper]{article}
\usepackage[margin=26mm,headheight=15pt]{geometry}
\usepackage{fontspec}
\setmainfont{TeX Gyre Pagella}
\usepackage{amsmath,amssymb,amsthm,mathtools,enumitem}
\usepackage{xcolor,hyperref,bookmark,fancyhdr}
\definecolor{linkblue}{HTML}{164B76}
\hypersetup{colorlinks=true,linkcolor=linkblue,urlcolor=linkblue,citecolor=linkblue,
 pdftitle={Deformation invariance of plurigenera in fiberwise Kähler families},
 pdfauthor={Tseng}}
\pagestyle{fancy}\fancyhf{}
\fancyhead[L]{\small Deformation invariance of plurigenera}
\fancyhead[R]{\small 11 November 2026}
\fancyfoot[C]{\thepage}
\setlength{\parindent}{1.3em}\setlength{\parskip}{3pt}
\setlist{itemsep=4pt,topsep=5pt}
\emergencystretch=2em
\newtheorem{theorem}{Theorem}[section]
\newtheorem{lemma}[theorem]{Lemma}
\newtheorem{proposition}[theorem]{Proposition}
\newtheorem{corollary}[theorem]{Corollary}
\theoremstyle{definition}\newtheorem{definition}[theorem]{Definition}
\theoremstyle{remark}\newtheorem*{remark}{Remark}
\newcommand{\Q}{\mathbb Q}\newcommand{\C}{\mathbb C}
\newcommand{\OO}{\mathcal O}
\newcommand{\cF}{\mathcal F}\newcommand{\cX}{\mathcal X}
\newcommand{\dd}{\mathrm d}
\newcommand{\Supp}{\operatorname{Supp}}\newcommand{\Sing}{\operatorname{Sing}}
\newcommand{\Spec}{\operatorname{Spec}}\newcommand{\Proj}{\operatorname{Proj}}
\newcommand{\per}{\operatorname{per}}\newcommand{\Tr}{\operatorname{Tr}}
\title{Deformation invariance of plurigenera\\in fiberwise Kähler families}
\author{Tseng}
\date{11 November 2026}
\begin{document}
\maketitle
\begin{abstract}
We prove deformation invariance of every plurigenus in a proper
surjective holomorphic submersion over a disk whose fibers are connected
compact Kähler manifolds, in arbitrary dimension. From OpenAI's
logarithmic subadditivity and good divisorial decomposition results,
we derive a finite
ordinary minimal model program starting on a chosen smooth fiber with
pseudoeffective canonical divisor. Hacon--Xie's deformation theorem
extends this program over the analytic family. Its semiample terminal
central fiber is covered by the classical extension argument of Levine
and Kollár--Mori. We also give a canonical extension theorem for
central fibers that are smooth in codimension two and admit
a compact Kähler modification. The proof combines Cohen--Macaulay
reflexive canonical powers, logarithmic Hodge theory, and a Leray
comparison to obtain a central period injection. Levine's trace formula
then annihilates the extension obstruction. Character summands recover
every pluricanonical degree, and coherence converts finite-order lifts
into analytic sections.
\end{abstract}
\tableofcontents
\section{Introduction and preliminaries}

Deformation invariance of plurigenera for compact complex surfaces is
due to Iitaka \cite[Theorem~III]{Iitaka1970}. In arbitrary dimension,
Siu established invariance for smooth projective families, first for
fibers of general type \cite[Main Theorem, p.~661]{Siu1998} and subsequently in general
\cite[Corollary~0.3]{Siu2002}. He also formulated the corresponding
conjecture for compact Kähler manifolds
\cite[Conjecture~0.4]{Siu2002}. We prove the following statement.

\begin{theorem}\label{thm:main}
Let $\Delta=\{t\in\C:|t|<1\}$, let $\cX$ be a complex manifold, and let
$\pi:\cX\to\Delta$ be a proper surjective holomorphic submersion.
Suppose that every fiber $X_t$ is connected, compact, and Kähler.
Then, for every integer $m\ge1$, the function
\[
 t\longmapsto P_m(X_t)
 =\dim_\C H^0(X_t,K_{X_t}^{\otimes m})
\]
is constant on $\Delta$.
\end{theorem}

Recent work has established invariance in several Kähler settings.
Deng, Ning, Wang, and Zhou \cite[Theorem~1.1]{DNWZ2026} prove
invariance when the total space is Kähler and every fiber has nef
canonical bundle. Chen, Rao, and Wang
\cite[Theorems~1.2--1.4]{CRW2026} obtain local invariance for smooth
proper Kähler morphisms from the nefness of the central canonical
bundle. Li, Păun, Wang, and Zhou \cite[Theorem~1.2]{LPWZ2026}
prove invariance in arbitrary dimension for smooth proper families
with Kähler total space, without a nefness assumption.
Liu and Wan \cite[Theorem~1.1]{LW2026} establish the corresponding
logarithmic extension theorem for rational boundaries with relative
simple normal crossings and coefficients in $[0,1)$, again with
Kähler total space.

The threefold case is also studied by Hacon, Li, and Rao
\cite[Theorem~6.5]{HLR2026}. Feng, Li, and Rao
\cite[Theorems~1.4 and~1.5]{FLR2026} prove invariance for
$\Q$-factorial Kähler families of threefolds with canonical central
fiber and for smooth Kähler families whose central fiber has maximal
Albanese dimension.
In their terminology, a Kähler family has Kähler total space.

Theorem~\ref{thm:main} establishes invariance in arbitrary dimension
under the fiberwise Kähler hypothesis.
A related local question, with nef canonical bundle on the central
fiber, is posed in \cite[Question~1.6]{CRW2026}.
The cases of fiber dimension at most two are classical and are recalled
in Section~\ref{sec:main-proof}. In higher dimension, the argument is
local on the base and is carried out separately for each pluricanonical
degree. Upper semicontinuity settles the case of a vanishing central
plurigenus. Otherwise, a nonzero pluricanonical section makes the
central canonical divisor pseudoeffective, and the proof proceeds
through three constructions.

First, OpenAI's logarithmic subadditivity theorem
\cite[Corollary~6.2]{OAI033} supplies the reduced simple normal crossing
case of logarithmic Iitaka subadditivity assumed in
\cite[Assumption~1.1]{OAI034}. We therefore use the good divisorial
decomposition theorem \cite[Theorem~2.13]{OAI034} and the ordinary and
relative model constructions of that paper.
Proposition~\ref{prop:finite-original}, proved in
Section~\ref{sec:finite-program}, derives from these constructions
a finite ordinary $K$-MMP beginning with the
chosen smooth fiber and ending at a terminal model with semiample
canonical divisor. Its proof uses the good decomposition in the ambient
dimension and adapts the fixed-part, small-model, and nef-interval
constructions of \cite[Proposition~3.8 and Lemmas~3.11--3.12]{OAI034},
together with the comparison of negative multiplicities in the proof of
\cite[Theorem~3.13]{OAI034}.

Second, the deformation step is supplied by Hacon and Xie
\cite[Theorem~7.7]{HX2026}. Their theorem extends prescribed central
MMP steps in proper flat analytic families over a smooth curve with Kähler canonical
central fiber. Its Kähler hypothesis concerns the central fiber;
in our application the boundary and generalized nef data are zero,
so the additional boundary conditions are automatic.
Proposition~\ref{prop:finite-family} gives the ordinary, relatively
projective construction explicitly, including the central canonical
algebra and the comparison of pluricanonical sections in every degree.
The proof follows Kollár's deformation framework
\cite[paragraph~8]{Kollar2021}, using Nakayama's analytic restriction
and deformation results \cite[VI.3.8 and VI.5.2]{Nakayama2004}.

Finally, the semiample terminal extension result needed at the endpoint
is classical: the proof of Kollár--Mori
\cite[Theorem~12.5.5]{KM1992} applies in Fujiki's class $\mathcal C$,
as recorded by Kollár \cite[comments following Conjecture~17]{Kollar2021}.
Theorem~\ref{thm:terminal-extension} extends this terminal
class-$\mathcal C$ formulation to canonical central fibers smooth in
codimension two and admitting a compact Kähler modification.
We prove this canonical statement using the cyclic-cover method of Levine
\cite[Proposition~1.6 and Corollaries~1.8 and~1.10]{Levine1983}
and Kollár--Mori. Cohen--Macaulayness identifies the relevant first
cohomology with that of the smooth locus.
On a compact Kähler resolution of the central cover, Deligne's
logarithmic Hodge-to-de Rham degeneration, in the formulation of
Liu, Rao, and Wan \cite[Theorem~3.2]{LRW2019}, gives an injection
into logarithmic de Rham cohomology. A Leray comparison and de Rham
naturality then imply that the central period map is injective.
Its infinitesimal counterpart annihilates the extension obstruction.
The character summands recover every lower pluricanonical degree, and
coherence converts finite-order lifts into analytic sections.
Birational invariance completes the proof in
Section~\ref{sec:main-proof}.

\subsection{Notation and conventions}

The letter $n$ denotes the dimension of the fibers and their
bimeromorphic models.
For a normal $\Q$-Gorenstein space, $\omega_X^{[m]}$ denotes the reflexive
canonical power. The notation $L\sim_\Q L'$ means that there is a positive
integer $r$ for which $rL$ and $rL'$ are holomorphic line bundles and are
isomorphic.
Semiampleness means global generation of a positive Cartier multiple.

A normal space is \emph{globally strongly $\Q$-factorial} if every global
coherent rank-one reflexive sheaf has an invertible reflexive power.
In particular, every global Weil divisor has a positive Cartier multiple.
Here the reflexive sheaves are taken on the whole space.
A bimeromorphic map is \emph{small} if it is an isomorphism in codimension one.
An ordinary \emph{$K$-minimal model program}, or \emph{$K$-MMP}, consists
of divisorial contractions and flips of $K$-negative extremal rays.
Braces $\{D\}$ denote the Bott--Chern
class of a divisor. In formulas involving classes, a rational line bundle
is identified with its first Chern class.
Canonical comparisons are understood through local meromorphic
canonical identifications; their discrepancy divisors are global.

Following Boucksom's analytic divisorial Zariski decomposition,
we write $N(\alpha)$ for the divisorial negative part
of a pseudoeffective Bott--Chern class.
On a smooth compact Kähler model, with Kähler form
$\omega$, the divisorial minimal multiplicities are
\[
 \nu_D(\alpha)=\lim_{\epsilon\downarrow0}
 \inf_{T\in\alpha,\ T\ge-\epsilon\omega}\nu_D(T),
 \qquad N(\alpha)=\sum_D\nu_D(\alpha)D.
\]
The support of this sum is finite
\cite[Theorem~3.12\textup{(i)}]{Boucksom}.
The class is \emph{modified nef} when
all its minimal multiplicities along prime divisors on the current space
vanish. On normal models, we use these notions on smooth resolutions and
keep track of all divisorial places. Comparisons between models are made
on a common smooth resolution.

\subsection{Birational results for compact Kähler spaces}

We use the following consequences of \cite{OAI033,OAI034}. The
statements concern compact Kähler spaces and the indicated proper or
projective morphisms.

\begin{enumerate}[label=\textbf{(B\arabic*)},leftmargin=*]
\item\label{in:decomp}
\textbf{Good canonical decomposition.}
If $X$ is smooth compact Kähler and $K_X$ is pseudoeffective, there is a
smooth compact Kähler modification $\mu:Y\to X$ such that
\begin{equation}\label{eq:good}
 \mu^*K_X\sim_\Q P+R,\qquad
 P\text{ is semiample},\quad R=N(\mu^*K_X)\ge0,
\end{equation}
where $R$ is a rational divisor. On every higher smooth modification $a$,
one has $N(a^*\mu^*K_X)=a^*R$.
This is \cite[Theorem~2.13 and Lemma~2.4]{OAI034}.

The morphism $\mu$ may be chosen projective. Indeed, apply analytic
flattening \cite{Hironaka1975} to the original modification. After a
sequence of blowups of $X$, its main transform is flat. It is proper of
relative dimension zero and degree one, hence finite flat of degree one.
The associated rank-one unital algebra is $\OO$, so the main transform
is an isomorphism. A projective resolution then gives a projective
modification of $X$ dominating the original one. Pullback preserves
\eqref{eq:good} and the asserted negative-part identities.

\item\label{in:steps}
\textbf{Ordinary steps and finiteness at positive scaling parameters.}
On a globally strongly $\Q$-factorial compact Kähler terminal space,
we use the ordinary scaling programs constructed in
\cite[Lemma~3.6]{OAI034}. Their steps are projective contractions
and flips of $K$-negative extremal rays. The resulting
spaces remain compact Kähler, terminal, and globally strongly
$\Q$-factorial. The rays are tested against all Bott--Chern classes.
The cone theorem recalled in \cite[Theorem~1.3]{HX2026} gives the
rational-curve length bound $2n$. Following
\cite[Section~3.1]{OAI034}, we use the uniform bound $N_n=2n+1$.

Starting with a Kähler scaling class, such a $K$-program can be run,
and only finitely many steps occur in every parameter interval
$[\epsilon,T]$ with $\epsilon>0$.
If $K$ is pseudoeffective and the program does not reach a nef model
after finitely many steps, its positive scaling parameters tend to zero.
The corresponding statements apply to $K+bP$, where $P$ is a nef
rational line bundle on the current model: if $mP$ is Cartier and
$b\in\Q$ satisfies $b>4nm$, the steps can be chosen
$K$-negative and $P$-trivial.

The effective comparisons are those of \cite[Lemma~3.7]{OAI034}.
The projective construction of the steps and finiteness at positive
scaling parameters are supplied by
\cite[Proposition~3.42 and Corollary~3.50]{OAI034}.

\item\label{in:descent}
\textbf{Fixed Cartier indices and transport of classes.}
If a Cartier line bundle $mP$ has degree zero on a contraction
as above, it descends as a Cartier line bundle to the target, with the
same integer $m$. Nefness and global generation of a fixed multiple
are preserved. A real Bott--Chern class of degree zero on the contracted
ray also descends to the contraction target.
See \cite[Lemmas~3.3 and~3.15]{OAI034}.

The fixed-index assertion can also be seen from relative klt
base-point-freeness: two consecutive sufficiently large powers descend,
and their quotient descends the original line bundle. The same Cartier
exponent is therefore available throughout the program.

\item\label{in:bigrelative}
\textbf{Big relative good models.}
Let $M$ be compact Kähler, terminal, and globally strongly
$\Q$-factorial, with semiample $K_M$. Let $g:M\to Z$ be its associated
morphism with connected fibers, where $Z$ is normal projective.
Let $A_M$ be modified nef and big, represented by globally nef
generalized data on a fixed smooth carrier. If
$L=K_M+tA_M$, for some $t>0$, defines a generalized klt pair, then $L$
has a nonextracting good log terminal model over $Z$.
The output is again globally strongly $\Q$-factorial and generalized
klt. Its comparison error is effective, exceptional over the output,
and strictly positive along each contracted prime divisor of the source.
The construction preserves the morphism to $Z$ and the forward
Bott--Chern traces.

Here $L$ is globally big, since $K_M$ is nef and $tA_M$ is big.
We use the proper-relative good-model assertion of
\cite[Proposition~3.2(ii) and Theorem~3.49]{OAI034}, with its globally
big construction retaining the chosen base \cite[Lemma~3.45]{OAI034}.
In Lemma~\ref{lem:nef-interval}, that base is the image of the
semiample morphism defined by $K_M$; properness of this morphism is sufficient.
The generalized klt adjoints occurring here satisfy the same uniform
rational-curve length bound $N_n$.

\item\label{in:projrelative}
\textbf{Projective birational relative programs.}
Given a projective bimeromorphic morphism from a smooth compact Kähler
space to a normal compact Kähler target, the ordinary relative
$K$-program can be run and terminated. Its terminal output remains
projective over that target and globally strongly $\Q$-factorial.
See \cite[Proposition~3.4]{OAI034}.
\end{enumerate}

\subsection{Divisorial negative parts}

We use exceptional translation on resolutions of normal Kähler spaces
in the form of \cite[Lemma~2.3]{OAI034}. The subtraction and fixed-section
properties are recalled in \cite[Lemmas~2.2 and~2.5]{OAI034}.
If $q:W\to T$ is a modification, $A$ is pseudoeffective, and $E\ge0$
is $q$-exceptional, then
\begin{equation}\label{eq:exceptional}
 N(q^*A+E)=N(q^*A)+E.
\end{equation}
If $0\le F\le N(A)$, then
\begin{equation}\label{eq:subtract}
 N(A-F)=N(A)-F.
\end{equation}
Nef classes have zero negative part, addition of a nef class does not
increase divisorial minimal multiplicities, and strict-transform
multiplicities are unchanged under modifications. The divisor of any
nonzero section of a line bundle contains the corresponding negative part.

We also use the following nef-direction form of the continuity argument
in \cite[Lemma~3.7]{OAI034}. If $A$ is pseudoeffective and $B$ is nef,
then, for every
fixed divisorial place $Q$,
\begin{equation}\label{eq:nef-limit}
 \nu_Q(A+\epsilon B)\longrightarrow\nu_Q(A)
 \qquad(\epsilon\downarrow0).
\end{equation}
To see this, work on a smooth model carrying $Q$, choose a Kähler class
$\Omega$, and choose $c$ sufficiently large that $c\Omega-B$ is nef.
Then
\[
 \nu_Q(A+c\epsilon\Omega)
 \le\nu_Q(A+\epsilon B)\le\nu_Q(A).
\]
The left side tends to the right side by the definition of minimal
multiplicities. This proves \eqref{eq:nef-limit}. By homogeneity, the
same observation applies to $(1-s)A+sB$.

Finally, an exceptional real divisor with zero Bott--Chern class is
zero. Both it and its negative are relatively nef, so the negativity
lemma applies to each sign. We shall use this observation to recover
exceptional-divisor identities from class comparisons.

\section{Finite minimal model programs on the original space}
\label{sec:finite-program}

We derive a finite ordinary $K_X$-MMP beginning with the given
manifold $X$ from the birational constructions of \cite{OAI034}.
This form of the conclusion supplies the individual contractions and
flips to be extended over the family in Section~\ref{sec:family-program}.
The Kodaira-dimension-zero case is contained in
\cite[Corollary~3.10]{OAI034}.

OpenAI's special-termination theorem
\cite[Theorem~3.13]{OAI034} asserts that the exceptional, flipping,
and flipped loci of a chosen dlt scaling eventually avoid the supports
of the boundary floors on the successive models. For zero boundary,
that condition places no restriction on
the remaining flips. The proposition below establishes finiteness of a
chosen entire ordinary $K$-program, using the good decomposition
\cite[Theorem~2.13]{OAI034} in the ambient dimension.

\begin{proposition}\label{prop:finite-original}
Let $X$ be a smooth compact Kähler manifold with $K_X$ pseudoeffective.
Then there is a finite ordinary $K_X$-MMP whose endpoint $X_{\min}$ is a
compact Kähler terminal space with semiample canonical divisor
$K_{X_{\min}}$.
\end{proposition}

The proof reduces a potentially infinite scaling to an ambient terminal
model $T$ with $N(K_T)=0$, and restarts with a Kähler scaling direction.
A small comparison model then makes the transforms of
$K_T+sA$ nef on one interval containing $s=0$. On a common resolution,
every divisorial negative multiplicity is consequently affine throughout
that interval. The comparison in the proof of
\cite[Theorem~3.13]{OAI034} shows that a nontrivial block of negative flips
at a scaling wall changes the slope of at least one such multiplicity.
This forces a restarted
program on $T$ to terminate; adjoining the finite initial segment gives
the required program on $X$.

\subsection{Contraction of the fixed part}

The following is the terminal, boundary-free case of OpenAI's fixed-part
contraction theorem \cite[Proposition~3.8]{OAI034}. We recall the argument
with the scaling comparison that will also be used in the final step.

\begin{lemma}\label{lem:contract-fixed}
Let $T$ be a compact Kähler terminal space that is globally strongly
$\Q$-factorial, and suppose that
\[
 K_T\sim_\Q P_T+D_T,\qquad P_T\text{ nef},\quad D_T\ge0,
\]
where the terms are rational line bundles or rational divisors, and
$N(p^*K_T)=p^*D_T$ on every smooth resolution $p$. Then there is a finite
$K_T$-negative, $P_T$-trivial ordinary MMP whose endpoint satisfies
$K\sim_\Q P$.
If $P_T$ is semiample, then the canonical divisor of the endpoint is semiample.
\end{lemma}

\begin{proof}
Choose $m$ such that $mP_T$ is Cartier, and then choose an integer
$b>\max\{4nm,mN_n\}$. Put $H=K_T+bP_T$ and choose a Kähler class $A$
with $H+A$ nef. By \ref{in:steps}, run the $H$-MMP with scaling of $A$,
and denote its successive forward traces by $A_i$.
An $H$-negative ray generated by $C$ is also
$K$-negative, and its generator may be chosen so that
$0<-K\cdot C\le N_n$.
If $P\cdot C>0$, the fixed Cartier index gives $P\cdot C\ge1/m$, whence
$H\cdot C\ge-N_n+b/m>0$, a contradiction. Thus every step is $P$-trivial.
Statement \ref{in:descent} preserves both the same index $m$ and nefness.

Let $T_i$ be the $i$th working model, with $T_0=T$, set $\lambda_{-1}=1$, and let
$\lambda_i$ be its next scaling threshold. Then $H_i+\lambda A_i$ is nef
for $\lambda\in[\lambda_i,\lambda_{i-1}]$; consecutive thresholds may
coincide. At a step with $\lambda_i>0$, let $E_i\ge0$ be the
$H$-comparison divisor on a common model.
The class at the threshold is pulled back from the contraction
base, so its two pullbacks agree. Consequently, for $0\le\lambda\le\lambda_i$,
\begin{equation}\label{eq:scaling-comparison}
 p_i^*(H_i+\lambda A_i)-q_i^*(H_{i+1}+\lambda A_{i+1})
   =\left(1-\frac{\lambda}{\lambda_i}\right)\{E_i\}.
\end{equation}
For a finite initial segment and a parameter $\lambda\ge0$ no larger than
any threshold crossed in that segment, summing these comparisons gives an
effective error divisor exceptional over the output.
If the corresponding class on the output is nef,
\eqref{eq:exceptional} identifies this error with the entire divisorial negative
part of the original class on the common model. We use this formula below for
$H=K+bP$, and later for $H=K$.

On the initial model and on every higher model, \eqref{eq:subtract} gives
\[
 H\sim_\Q(1+b)K_T-bD_T,
 \qquad N(H)=(1+b)D_T-bD_T=D_T.
\]
The descent of the line bundle $P$ and the canonical comparison preserve
this identity along the MMP: on a common model one has
$p^*D=q^*D_i+E$, and \eqref{eq:exceptional} gives the same formula for the
negative part. The prime divisor contracted by a divisorial step is covered by
contracted curves. Since $P$ has degree zero on those curves,
$D\cdot C=K\cdot C<0$. An effective rational Cartier divisor has nonnegative
degree on every curve not contained in its support, as follows after clearing
denominators. Hence the contracted prime divisor
belongs to $\Supp D$.
If the MMP terminates at a model on which $H_i$ is nef, then
$D_i=N(H_i)=0$, so $K_i\sim_\Q P_i$.

It remains to exclude an infinite MMP. In that case its positive scaling
parameters tend to zero. By \eqref{eq:nef-limit}, each component of the original
$D_T$ has positive coefficient in $N(H+\lambda A)$ for all sufficiently small
positive $\lambda$. On a nef model for such a parameter, the negative part is
entirely computed by the effective exceptional comparison divisor. If the
component survived as a strict prime divisor on that model, its coefficient
in the exceptional comparison would be zero, a contradiction. Thus each
component of $D_T$ is contracted at some finite stage. There are only finitely
many components. After all these stages, $K\sim_\Q P$, so $K$ is nef
and the MMP terminates. The Cartier descents preserve a fixed globally generated multiple,
so they also preserve semiampleness.
\end{proof}

\subsection{Canonical decompositions on intermediate models}

Fix \eqref{eq:good}. Let $T$ be a terminal model obtained from $X$ by a finite
ordinary MMP. Choose a smooth common model $W$ dominating $Y$, $X$ and $T$,
and write $r:W\to X$ and $s:W\to T$ for the resulting morphisms.
The canonical meromorphic comparisons give
\[
 K_W=r^*K_X+E_X=s^*K_T+E_T,
 \qquad E_X,E_T\ge0,
\]
where the error divisors are exceptional for the respective morphisms.
The finite $K$-negative program has the canonical comparison
$r^*K_X=s^*K_T+G$, where $G\ge0$ is $s$-exceptional.
If $\Theta$ is a positive current in $c_1(K_X)$, then $s_*r^*\Theta$
is a positive current in $c_1(K_T)$, since $s_*[G]=0$.
Hence $K_T$ is pseudoeffective. Pulling back \eqref{eq:good} and applying
\eqref{eq:exceptional}, we obtain
\[
 K_W\sim_\Q P_W+N_W,
 \qquad P_W\text{ semiample},\quad N_W=N(K_W).
\]
A second application of \eqref{eq:exceptional} gives
\begin{equation}\label{eq:model-good}
 R_T:=N_W-E_T=N(s^*K_T)\ge0,
 \qquad s^*K_T\sim_\Q P_W+R_T.
\end{equation}
Thus the initial good decomposition induces \eqref{eq:model-good} on every
finite intermediate model.

\subsection{A small semiample comparison model}

We adapt the normalized-graph construction of
\cite[Lemma~3.11]{OAI034} to canonical divisors on terminal spaces.

\begin{lemma}\label{lem:small-semiample}
Suppose that the model $T$ above satisfies $N(K_T)=0$, with vanishing at every
prime divisor on $T$. Then there is a globally strongly $\Q$-factorial
compact Kähler terminal model $M$ that is small bimeromorphic to $T$
and has semiample canonical divisor $K_M$.
The model $M$ will serve as a comparison model for an ordinary $K_T$-MMP.
\end{lemma}

\begin{proof}
The divisor $R_T$ in \eqref{eq:model-good} is now $s$-exceptional. Choose $m$
sufficiently divisible that $mP_W$ is globally generated and $mR_T$ is integral.
The fixed-part property gives
\[
 H^0(T,mK_T)=H^0(W,ms^*K_T)
 =s_{mR_T}\,H^0(W,mP_W).
\]
Let $\Gamma\to T$ be the normalization of the graph of this complete
pluricanonical map. This morphism is projective, and $\Gamma$ carries a
semiample rational line bundle $P_\Gamma$ whose pullback to a sufficiently high
model over $W$ is precisely $P_W$.
Choose a resolution $W_0\to\Gamma$ with $W_0$ smooth and the morphism
projective, and apply statement
\ref{in:projrelative} to run an ordinary relative $K$-MMP over $\Gamma$.
Write its endpoint as $g:Y_0\to\Gamma$, and let $h:Y_0\to T$ be the composite.
The space $Y_0$ is terminal. Since $T$ is terminal,
\[
 K_{Y_0}=h^*K_T+F,
\]
where $F\ge0$ is $h$-exceptional and has strictly positive coefficient at every
$h$-exceptional prime divisor.

Set $P_0=g^*P_\Gamma$. Pushing down \eqref{eq:model-good} from a common model
gives
\[
 h^*K_T\sim_\Q P_0+R_0,
 \qquad R_0\ge0\text{ is an }h\text{-exceptional rational divisor}.
\]
Since $Y_0$ is globally strongly $\Q$-factorial, $R_0$ is rational Cartier, so
the following pullbacks are defined. On a common model
$a:W'\to W$, $v:W'\to Y_0$, the divisor $a^*R_T-v^*R_0$ is $v$-exceptional
and rationally linearly equivalent to zero. Applying negativity with both
signs shows that it vanishes. Therefore
\[
 K_{Y_0}\sim_\Q P_0+(R_0+F),\qquad
 N(v^*K_{Y_0})=v^*(R_0+F).
\]
The second identity also follows directly by applying \eqref{eq:exceptional}
to $hv$.

Lemma \ref{lem:contract-fixed} contracts $R_0+F$. Its support contains every
$h$-exceptional prime divisor and contains no prime divisor coming from $T$.
Moreover, every divisor contracted by the MMP belongs to this support.
The resulting model $M$ is therefore small bimeromorphic to $T$, terminal and
Kähler, with $K_M\sim_\Q P_M$ semiample. It is globally strongly
$\Q$-factorial by statements~\ref{in:projrelative} and~\ref{in:steps}.
\end{proof}

\subsection{A common nef model for an interval}

The construction of \cite[Lemma~3.12]{OAI034} produces one model that
is nef for an entire interval of adjoints. We record its exceptional
comparison and the resulting affine formula for the negative part,
as used in the proof of \cite[Theorem~3.13]{OAI034}.

\begin{lemma}\label{lem:nef-interval}
In the setting of Lemma \ref{lem:small-semiample}, let $A$ be a Kähler class
on $T$. There are a globally strongly $\Q$-factorial compact Kähler
terminal model $V$ small bimeromorphic to $T$, a forward trace $A_V$
of this class, and $\delta>0$ such that
\[
 K_V+sA_V\text{ is nef}\qquad(0\le s\le\delta).
\]
On a common model $p:U\to T$, $q:U\to V$, there is an effective real divisor
$F(s)$, exceptional for both morphisms and with coefficients affine in $s$,
such that
\begin{equation}\label{eq:affine-N}
 \begin{aligned}
 p^*(K_T+sA)&=q^*(K_V+sA_V)+\{F(s)\},\\
 N\bigl(p^*(K_T+sA)\bigr)&=F(s)
 \end{aligned}
\end{equation}
throughout $[0,\delta]$. This affine dependence persists on every higher smooth
model.
\end{lemma}

\begin{proof}
First transport $h^*A$ to $M$ along the finite construction of
Lemma \ref{lem:small-semiample}. At
each $K_i$-negative step, choose a generating curve $C_i$ and set
$c_i=(\alpha_i\cdot C_i)/(K_i\cdot C_i)$.
The class $\alpha_i-c_iK_i$ vanishes on the entire analytic numerical ray.
By statement \ref{in:descent}, it descends to the contraction base. On the next
model, add back $c_iK_{i+1}$. The difference on a common model is $c_i$ times
the canonical comparison divisor. Summing finitely many such
comparisons produces $A_M$. On a common model $p:U\to T$, $q:U\to M$,
\begin{equation}\label{eq:A-transport}
 q^*A_M-p^*A=\{F_A\},
\end{equation}
where $F_A$ is exceptional for both morphisms. The divisor $-F_A$ is $q$-nef,
since it has the same degree as $p^*A$ on every $q$-contracted curve.
Negativity gives $F_A\ge0$. Applying \eqref{eq:exceptional} to $p$ yields
$N(q^*A_M)=F_A$, so $A_M$ is modified nef. It is also big, since $p^*A$ is big.

Set $P_M=K_M$. On the fixed common carrier, take $tp^*A$ as the globally nef
b-data, and take zero ordinary boundary on $M$. This gives the generalized
adjoint $L=P_M+tA_M$. On a fixed log resolution, its boundary is the ordinary
crepant boundary plus $tF_A$. Since $M$ is terminal, all coefficients remain
strictly less than one for a sufficiently small fixed $t>0$, so the pair is
generalized klt. The total boundary-plus-nef trace is globally modified big,
and $L$ is globally big and modified nef.

Write $mP_M=g_M^*H$, where $g_M:M\to Z$ has connected fibers, $Z$ is normal
projective, and $H$ is ample Cartier. Applying statement~\ref{in:bigrelative}
gives a globally strongly $\Q$-factorial good model
$M\dashrightarrow V$ over $Z$. Set $P_V=g_V^*H/m$.
This map is small. Indeed, on a common model one has
$a^*L=c^*L_V+E$, with $E\ge0$ and with strictly positive coefficient at every
original prime divisor that is contracted. The class $L_V$ is a pseudoeffective
forward trace, so \eqref{eq:exceptional} gives $N(a^*L)\ge E$.
Since $L$ is modified nef, its negative part has coefficient zero at every
original prime divisor. This excludes all divisorial contractions;
nonextraction excludes divisors in the other direction.

The map is over $Z$ and small, so $K_V\sim_\Q P_V$.
On a common model, the difference of the two canonical pullbacks is
exceptional and has zero class, hence vanishes. Thus $V$ is still terminal.
Define $A_V=(L_V-P_V)/t$.
Choose $b>mN_n$. Then $L_V+bP_V$ is globally nef. Otherwise a negative ray
has a generating curve $C$ with $0<-L_V\cdot C\le N_n$.
If $P_V\cdot C=0$, then $C$ is vertical over $Z$, contrary to the relative
nefness of $L_V$. If $P_V\cdot C>0$, then $P_V\cdot C\ge1/m$, again
contradicting $b>mN_n$.

Both $P_V$ and $(1+b)P_V+tA_V$ are nef. Convexity therefore gives
\[
 P_V+sA_V\text{ is nef},\qquad 0\le s\le\delta=\frac{t}{1+b}.
\]
Pull all comparisons to a common smooth model of $T$ and $V$. Since the
bimeromorphic map is small, the difference $F(s)$ is exceptional for both
morphisms, and the linear transport makes its coefficients affine in $s$.
The divisor $-F(s)$ is $p$-nef because $q^*(K_V+sA_V)$ is nef. Negativity gives
$F(s)\ge0$. Applying \eqref{eq:exceptional} to $q$ proves
\eqref{eq:affine-N}. On a higher model, the pullback of the nef class remains
nef and the pullback of the error remains exceptional. Hence the entire
statement persists.
\end{proof}

\subsection{Termination of the MMP}

\begin{proof}[Proof of Proposition \ref{prop:finite-original}]
Starting from the original smooth manifold $X$, run the ordinary $K_X$-MMP
with scaling supplied by statement \ref{in:steps}. If it reaches a nef endpoint
after finitely many steps, proceed to the final paragraph of the proof.
Suppose that it is infinite. Every original prime divisor that is contracted
has positive coefficient in $N(K_X)$: its coefficient in the canonical
comparison at the relevant finite stage is strictly positive, while the
canonical class on the other side remains pseudoeffective. Formula
\eqref{eq:exceptional} gives the asserted lower bound. The MMP is
nonextracting, and $N(K_X)$ has finite support. Hence only finitely many
divisorial steps occur.

After the last divisorial step, pass also beyond the current finite block of
steps having the same threshold, and denote the resulting model by $T$.
All subsequent steps are flips, and $T$ has a nonempty interval of positive
parameters on which it is a nef model. Denote the trace of the original
scaling direction by $H$; it need not be nef. Choose $\lambda_0>0$ in that
interval so that $B=K_T+\lambda_0H$ is nef.
For subsequent scaling parameters $\lambda>0$ tending to zero, the small nef
models give the exceptional comparisons for $K_T+\lambda H$. Its negative
part therefore has coefficient zero at every prime divisor on $T$.
Using
\[
 K_T+\lambda H
 =\left(1-\frac{\lambda}{\lambda_0}\right)K_T
  +\frac{\lambda}{\lambda_0}B
\]
and \eqref{eq:nef-limit}, we obtain $N(K_T)=0$.
The limit argument uses the nef class $B=K_T+\lambda_0H$.

Restart the ordinary $K_T$-MMP with scaling of a Kähler class $A$ on $T$.
Since the negative part of $K_T$ is zero, the new MMP contracts no divisor:
the canonical comparison argument above would force a strictly positive
generic multiplicity $\nu$ at any contracted prime divisor. Thus all its
nontrivial steps are small. Lemma \ref{lem:nef-interval} supplies a fixed
comparison model $V$ and an interval $[0,\delta]$.

If the new MMP were infinite, its positive parameters would tend to zero.
Every interval bounded away from zero contains only finitely many steps, so
there would be infinitely many distinct positive thresholds. Combine the
finitely many consecutive steps at each threshold $s_i$ into a single block.
On a common resolution of the models immediately before and after this block,
let $D_i\ge0$ be the canonical comparison divisor. Strict $K$-negativity of
the steps gives $D_i\ne0$. At the wall $s=s_i$, the pullbacks of the scaled
classes agree. Thus, for every $s$,
\begin{equation}\label{eq:wall}
 p_i^*(K_i+sA_i)-q_i^*(K_{i+1}+sA_{i+1})
   =\left(1-\frac{s}{s_i}\right)\{D_i\}.
\end{equation}
On each side of the wall, consecutive distinct thresholds bound a nonempty
open interval. The corresponding nef models compute the entire negative part
of the original class $K_T+sA$.
On a resolution simultaneously dominating these models and the fixed model
$V$, choose a prime divisor whose coefficient $d_i$ in $D_i$ is positive.
Equation \eqref{eq:wall} shows that the slopes of the affine formulas for
this divisorial multiplicity on the two sides of the wall differ by
$-d_i/s_i\ne0$. Choose $0<s_i<\delta$. This contradicts the \emph{single}
affine formula in \eqref{eq:affine-N} on the whole interval $[0,\delta]$.
The new MMP is therefore finite. Concatenating it with the finite initial
segment of the original MMP ending at $T$ gives a finite MMP starting
from the original $X$.

Let $X_{\min}$ be its endpoint, with $K_{X_{\min}}$ nef.
On a resolution, \eqref{eq:model-good} gives
$s^*K_{X_{\min}}\sim_\Q P_W+N(s^*K_{X_{\min}})$.
Nefness makes the negative part vanish. Choose a positive integer $q$ such
that $qK_{X_{\min}}$ is Cartier and
$s^*\OO_{X_{\min}}(qK_{X_{\min}})$ is globally generated.
This is possible by taking a common multiple of the Cartier index downstairs
and a generating multiple upstairs. Since $s_*\OO_W=\OO_{X_{\min}}$,
the projection formula identifies their spaces of global sections.
A base point downstairs would make its entire nonempty inverse image a base
locus upstairs, a contradiction. Thus $K_{X_{\min}}$ is semiample.
\end{proof}

\begin{remark}
Strict $K$-negativity is used in \eqref{eq:wall} to obtain $D_i\ne0$.
The deformation argument below applies to the resulting finite ordinary
$K_X$-MMP beginning with $X$.
\end{remark}

\section{Deformation of finite minimal model programs}
\label{sec:family-program}

Hacon and Xie \cite[Theorem~7.7]{HX2026} prove that prescribed
central MMP steps extend in proper flat analytic families over a smooth
curve with Kähler canonical central fiber.
Applied with zero boundary and zero generalized nef
data, their theorem supplies the deformation step needed after
Proposition~\ref{prop:finite-original}.
We give the ordinary, relatively projective construction in
Proposition~\ref{prop:finite-family}, identifying the central flip
algebra in every degree and the pluricanonical section spaces on
nearby fibers.

The argument follows the deformation framework of Kollár
\cite[paragraph~8]{Kollar2021}. Its form for generalized pairs with
projective central fibers is developed by Hacon, Li, and Rao
\cite[Proposition~4.16]{HLR2026}, with the Kähler threefold case
noted in \cite[Remark~4.17]{HLR2026}; Hacon--Xie treat
Kähler central fibers in arbitrary dimension. See also
\cite[Theorem~2.28]{FLR2026}.

\subsection{Birational invariance of reflexive pluricanonical sections}

We begin with the standard birational invariance statement in the
sheaf-theoretic form used to compare the fibers of the successive models.

\begin{lemma}\label{lem:canonical-birational}
Let $Z$ be a normal complex space with canonical singularities, and let
$p:W\to Z$ be a proper resolution with $W$ smooth. Under the usual
identification of canonical differentials, one has
\begin{equation}\label{eq:canonical-resolution-sections}
 p_*\omega_W^{\otimes m}=\omega_Z^{[m]}
\end{equation}
for every integer $m\ge1$. Consequently, if $Z_1$ and $Z_2$ are
bimeromorphic compact normal spaces with canonical singularities, then
\begin{equation}\label{eq:canonical-birational-H0}
 H^0(Z_1,\omega_{Z_1}^{[m]})
 \simeq H^0(Z_2,\omega_{Z_2}^{[m]})
 \qquad(m\ge1).
\end{equation}
The identities hold for the reflexive canonical powers in every degree.
\end{lemma}

\begin{proof}
The asserted sheaf identity is local on $Z$. On each sufficiently small
open subset of $Z$, choose a meromorphic canonical differential and use
its pullback to define compatible local canonical divisors on $W$.
Write
\[
 K_W=p^*K_Z+E,\qquad E\ge0,
\]
where $E$ is a $p$-exceptional rational divisor. Let $\eta$ denote the
chosen meromorphic top differential. A local section $s$ of the
reflexive $m$th canonical power can be written as
$s=u\eta^{\otimes m}$, with
\[
 D_s:=\operatorname{div}(u)+mK_Z\ge0.
\]
Since $K_Z$ is rational Cartier, $D_s$ is an effective rational Cartier
divisor. Its pullback is therefore effective. Viewed as a meromorphic
$m$-canonical differential on $W$, the pullback of $s$ has divisor
\[
 \operatorname{div}(p^*u)+mK_W=p^*D_s+mE\ge0,
\]
and is thus holomorphic. Conversely, a holomorphic $m$-canonical
differential on $W$ restricts to the common dense open set. It has no
pole at any codimension-one point of $Z$, and hence extends uniquely
as a section of the reflexive canonical power. This proves
\eqref{eq:canonical-resolution-sections} on every open subset.
Applying it on a common smooth resolution of $Z_1$ and $Z_2$ and taking
global sections gives \eqref{eq:canonical-birational-H0}.
\end{proof}

\subsection{Deformation of a prescribed finite program}

\begin{proposition}\label{prop:finite-family}
Let $f:\mathfrak X\to D$ be a proper flat analytic family over a disk
$D$ containing the origin. Suppose that $X_0$ is connected, normal,
compact, and has canonical singularities. Let
\begin{equation}\label{eq:prescribed-central-program}
 X_0=X_0^{(0)}\dashrightarrow X_0^{(1)}
 \dashrightarrow\cdots\dashrightarrow X_0^{(r)}
\end{equation}
be a prescribed finite ordinary $K$-MMP. Assume that every
central contraction
$\varphi_0^{(i)}:X_0^{(i)}\to Z_0^{(i)}$ is projective with connected
fibers and that $-K_{X_0^{(i)}}$ is relatively ample. Assume also that
every central output is normal with canonical singularities and, in
the flipping case, that the prescribed morphism
$X_0^{(i+1)}\to Z_0^{(i)}$ is the corresponding flip.

After shrinking $D$, there are proper flat analytic families
$\mathfrak X^{(i)}\to D$ whose central fibers are precisely the spaces
in \eqref{eq:prescribed-central-program}, and the prescribed steps
extend to bimeromorphic diagrams between these families. Near the
central fibers, their total spaces are normal, canonical, and
$\Q$-Gorenstein. Every nearby fiber is normal and canonical, and is
bimeromorphic to the corresponding original fiber $X_t$. In
particular, for every $m\ge1$ and all sufficiently small $t$,
\begin{equation}\label{eq:family-birational-plurigenera}
 H^0(X_t,\omega_{X_t}^{[m]})
 \simeq H^0(X_t^{(r)},\omega_{X_t^{(r)}}^{[m]}).
\end{equation}
\end{proposition}

We first extend the contraction underlying a step. For a flip, we identify
the central restriction of the relative canonical algebra, lift finite
generation from the prescribed central flip, and take relative Proj.
The following local and relative analytic results supply the inputs for
these constructions.

\begin{enumerate}[label=\textbf{(A\arabic*)},leftmargin=*]
\item\label{an:canonical-deform}
\textbf{Analytic deformations of canonical singularities.}
If the central fiber of a flat analytic family over a smooth disk has a canonical
singularity at a point, then the relative canonical divisor is rational
Cartier near that point and all nearby fibers have canonical
singularities there. The total space is canonical,
and the exceptional discrepancies of the pair $(\mathfrak X,X_0)$
are nonnegative. The relative $\Q$-Gorenstein property and deformation
statements are recorded in \cite[Lemma~10]{Kollar2021}; deformation of
canonical singularities is given by \cite[VI.5.2]{Nakayama2004}.
The total-space and pair statements over a smooth disk are recorded in
\cite[Remark~10.3]{Kollar2021}. Finitely many such
neighborhoods cover a compact central fiber. Properness then permits
a further shrinking of the disk so that their union contains the
entire restricted family.

\item\label{an:contraction-deform}
\textbf{Deformation of a contraction.}
Let $\varphi_0:U_0\to V_0$ be a proper contraction between compact
normal complex spaces. If
$(\varphi_0)_*\OO_{U_0}=\OO_{V_0}$ and
$R^1(\varphi_0)_*\OO_{U_0}=0$, then the contraction extends over every
given flat analytic deformation $\mathfrak U\to D$ of $U_0$:
after shrinking $D$, there is a flat analytic deformation
$\mathfrak V\to D$ of $V_0$ and a proper morphism
$\mathfrak U\to\mathfrak V$ extending $\varphi_0$.
This is the proper case of Kollár--Mori's deformation theorem
\cite[Proposition~11.4, part~(11.4.1)]{KM1992}.

\item\label{an:relative-restriction}
\textbf{Pluricanonical restriction for a projective bimeromorphic morphism.}
Let $g:Y\to Z$ be a projective bimeromorphic analytic morphism, with
$Y$ smooth and $Z$ normal. If $D_Y\subset Y$ is a smooth divisor
not contracted by $g$, then
\begin{equation}\label{eq:nakayama-relative-restriction}
 g_*\omega_Y^{\otimes m}(mD_Y)
 \twoheadrightarrow(g|_{D_Y})_*\omega_{D_Y}^{\otimes m}
 \qquad(m\ge1).
\end{equation}
This special case of Nakayama's relative restriction theorem
\cite[VI.3.8]{Nakayama2004} is stated in
\cite[Theorem~9]{Kollar2021}. Projectivity here is relative to $Z$.

\item\label{an:finite-generation-lift}
\textbf{Lifting graded finite generation near a central compact set.}
Let $g:Y\to Z$ be a proper morphism of normal complex spaces, let
$L$ be a line bundle on $Y$, and let $W\subset Z$ be a closed
analytic subspace. If
\[
 \OO_W\otimes_{\OO_Z}\bigoplus_{m\ge0}g_*L^{\otimes m}
\]
is finitely generated, then every compact subset of $W$ has a
neighborhood on which the original graded algebra is finitely
generated \cite[Lemma~11]{Kollar2021}. In the application below,
finite generation is supplied by the prescribed central flip.
\end{enumerate}

\subsection{Relative vanishing and extension of a contraction}

Consider one step and suppress its index. The central source $X_0$
is canonical, hence klt, and $-K_{X_0}$ is $\varphi_0$-ample.
Fujino's vanishing theorem for projective analytic morphisms
\cite[Theorem~1.1(ii)]{FujinoVanishing2025} gives
\begin{equation}\label{eq:contraction-R1-vanishing}
 R^j(\varphi_0)_*\OO_{X_0}=0\qquad(j>0).
\end{equation}
To apply it, take a projective log resolution $b:W\to X_0$ and write
\[
 K_W=b^*K_{X_0}+E,\quad E\ge0,\qquad
 A=\lceil E\rceil,\quad B=A-E.
\]
Then $B$ is an SNC boundary with coefficients in $[0,1)$, and one
has the identity of rational line bundles
\[
 A-(K_W+B)=b^*(-K_{X_0}).
\]
Apply the cited vanishing theorem to
$b:W\to X_0$, $\varphi_0:X_0\to Z_0$,
$L=\OO_W(A)$, $H=-K_{X_0}$, and direct-image degree $q=0$.
It gives $R^j(\varphi_0)_*b_*\OO_W(A)=0$.
Since $A$ is effective and $b$-exceptional, normality implies
$b_*\OO_W(A)=\OO_{X_0}$, proving
\eqref{eq:contraction-R1-vanishing}.

The central contraction satisfies
$(\varphi_0)_*\OO_{X_0}=\OO_{Z_0}$. By
\ref{an:contraction-deform}, after shrinking the disk it extends to
\[
 \varphi:\mathfrak X\longrightarrow\mathfrak Z,
 \qquad h:\mathfrak Z\longrightarrow D,
 \qquad f=h\varphi,
\]
where $\mathfrak Z\to D$ is proper and flat with central fiber
$Z_0$, and $\varphi$ is a proper contraction with connected
fibers. Openness of normality in a flat deformation permits a
shrinking for which $\mathfrak Z$ and its fibers are normal.
The stronger canonical property for the source follows from
\ref{an:canonical-deform}.

The same result makes $K_{\mathfrak X/D}$ rational Cartier.
Compactness of the central fiber allows us to choose one integer
$a>0$ such that $aK_{\mathfrak X/D}$ is Cartier near the central
fiber and $aK_{X_0}$ is Cartier. Adjunction on the smooth locus of
the central fiber, followed by extension on the normal fiber, gives
\[
 \OO_{\mathfrak X}(aK_{\mathfrak X/D})|_{X_0}
 \simeq\OO_{X_0}(aK_{X_0}).
\]
Thus $\OO_{\mathfrak X}(-aK_{\mathfrak X/D})$ restricts to a
$\varphi_0$-ample line bundle. Relative ampleness for a proper
analytic morphism is open on the target. Locally, this can be
verified using relative Serre vanishing for sufficiently high powers,
the evaluation maps, and the resulting projective embeddings; finitely
many target neighborhoods suffice along $Z_0$. Consequently the line
bundle is $\varphi$-ample over some neighborhood of $Z_0$.
Since $h$ is proper, we may shrink $D$ so that the entire restricted
target lies in this neighborhood. It follows that $\varphi$ is
projective and $-K_{\mathfrak X/D}$ is relatively ample.

\subsection{Bimeromorphicity on every nearby fiber}

Choose $z\in Z_0$ such that the central contraction is an
isomorphism near $z$, and such that both $z$ and its unique
preimage $x$ are smooth points of their central fibers. Such
points exist because $\varphi_0$ is bimeromorphic. A flat
morphism is smooth at a smooth point of its fiber, so the two
total spaces are smooth near $x$ and $z$. The vertical
differential of $\varphi$ is invertible, and its induced map on
the base is the identity. Its total differential is therefore
invertible. The inverse function theorem gives a local
isomorphism between neighborhoods of $x$ and $z$.

After shrinking, this is an isomorphism on the \emph{full preimage}
of a target neighborhood. Indeed, choose a sufficiently small source
neighborhood $U$ on which the local inverse is defined.
The set $\mathfrak X\setminus U$ is closed, and properness of
$\varphi$ makes its image closed. That image does not contain
$z$, since $\varphi^{-1}(z)=\{x\}$. Removing this image from
the target neighborhood and taking a smaller open subset $V$ gives
\begin{equation}\label{eq:family-common-open}
 \varphi^{-1}(V)\xrightarrow{\sim}V.
\end{equation}
Since $h$ is a submersion at $z$, the image $h(V)$ contains a
neighborhood of the origin. After a further shrinking,
\eqref{eq:family-common-open} meets every fiber.

The central spaces are connected, compact, and normal, and hence
have $h^0(\OO)=1$. Upper semicontinuity for proper maps implies that nearby
fibers also have $h^0(\OO)=1$, so they are connected.
Distinct irreducible components of a normal space are disjoint;
the nearby fibers are therefore irreducible. Each
$\varphi_t:X_t\to Z_t$ is proper and surjective and contains the
nonempty open isomorphism supplied by
\eqref{eq:family-common-open}. It is thus bimeromorphic.
The total morphism $\varphi$ is bimeromorphic as well.

If $\varphi_0$ is a divisorial contraction, the prescribed central
MMP gives canonical singularities on $Z_0$. Apply
\ref{an:canonical-deform} to $h:\mathfrak Z\to D$ and set
$\mathfrak X^+=\mathfrak Z$. This extends the step, with all the
asserted properties of the nearby total space, fibers, and relative
canonical divisor.

\subsection{The flip algebra and central base change}

Suppose now that $\varphi_0$ is small. After shrinking the disk,
the total morphism $\varphi$ is also small. If a horizontal
exceptional prime $F\subset\mathfrak X$ existed, properness
would force it to meet the central fiber. Its intersection with
$X_0$ would contain a divisor of $X_0$. Since both $F$ and
$\varphi(F)$ dominate the disk, $\varphi(F)\cap Z_0$ would have dimension at most
$\dim X_0-2$, contradicting the smallness of $\varphi_0$.
A vertical prime divisor must be an entire irreducible fiber.
Since the fiber maps are bimeromorphic, such a divisor cannot be
exceptional either. Components outside the central neighborhood
can be excluded by shrinking the disk.

Take a projective log resolution $\tau:Y\to\mathfrak X$ such
that the strict transform $D_Y$ of $X_0$ is smooth. Set
\[
 g=\varphi\tau:Y\to\mathfrak Z,\qquad L=\omega_Y(D_Y).
\]
Then $g$ is projective and bimeromorphic, and $D_Y$ is not
contracted by $g$. The assertion about the central pair in
\ref{an:canonical-deform} gives the canonical comparison
\begin{equation}\label{eq:family-canonical-pair-comparison}
 K_Y+D_Y=\tau^*(K_{\mathfrak X}+X_0)+E,
 \qquad E\ge0\text{ is }\tau\text{-exceptional}.
\end{equation}
The local differential argument of
Lemma~\ref{lem:canonical-birational}, now allowing poles of order
at most $m$ along $X_0$, gives
\[
 \tau_*L^{\otimes m}=
 \omega_{\mathfrak X}^{[m]}(mX_0)
 \qquad(m\ge0).
\]
Since $\varphi$ is small, the reflexive sheaves on both sides are
determined by their restrictions to the same open set containing
all codimension-one points. Hence
\begin{equation}\label{eq:family-canonical-algebra}
 \mathcal B_m:=g_*L^{\otimes m}
 =\omega_{\mathfrak Z}^{[m]}(mZ_0),
 \qquad \mathcal B=\bigoplus_{m\ge0}\mathcal B_m.
\end{equation}
Here $\omega_{\mathfrak Z}^{[m]}(mZ_0)$ denotes the reflexive
divisorial sheaf on the normal space $\mathfrak Z$.

Apply \ref{an:relative-restriction} to $g$ and $D_Y$.
The morphism $D_Y\to X_0$ is a smooth resolution. By
Lemma~\ref{lem:canonical-birational} and the smallness of
$\varphi_0$,
\[
 (g|_{D_Y})_*\omega_{D_Y}^{\otimes m}
 =(\varphi_0)_*\omega_{X_0}^{[m]}
 =\omega_{Z_0}^{[m]}.
\]
We therefore obtain surjections
\begin{equation}\label{eq:family-algebra-restriction}
 \mathcal B_m\twoheadrightarrow\omega_{Z_0}^{[m]}
\end{equation}
for every $m\ge1$, compatible with multiplication. In degree zero
the map is $\OO_{\mathfrak Z}\twoheadrightarrow\OO_{Z_0}$.

Let $t$ be the disk coordinate. The kernel of
\eqref{eq:family-algebra-restriction} is exactly $t\mathcal B_m$.
At a general smooth point of the central prime divisor $Z_0$,
the total family is smooth and the restriction map is ordinary
restriction under adjunction. Thus a local section $s$ with zero restriction
is divisible by $t$ there. At the general point of any other prime
divisor, $t$ is a unit. The meromorphic section $s/t$ of $\mathcal B_m$
is therefore holomorphic at every codimension-one point.
Since $\mathcal B_m$ is reflexive and $\mathfrak Z$ is normal,
$s/t$ extends uniquely across
the remaining subset of codimension at least two. The reverse
inclusion is immediate. Consequently one has an identity of
graded algebras
\begin{equation}\label{eq:family-algebra-base-change}
 \mathcal B/t\mathcal B
 =\bigoplus_{m\ge0}\omega_{Z_0}^{[m]}.
\end{equation}

The right-hand algebra is finitely generated, and its relative
Proj is precisely the \emph{prescribed} central flip $X_0^+$.
Indeed, writing $\varphi_0^+:X_0^+\to Z_0$ for the prescribed flip,
this morphism is small and $K_{X_0^+}$ is rational
Cartier and relatively ample, so its relative canonical algebra
agrees with the right-hand side in every degree. Choose $a>0$ such that
$aK_{X_0^+}$ is Cartier and relatively ample. This multiple gives
finite generation of the $a$th Veronese subalgebra.
For each $0\le r<a$, relative Serre finiteness makes
\[
 \bigoplus_{q\ge0}(\varphi_0^+)_*
 \bigl(\omega_{X_0^+}^{[r]}\otimes
 \OO_{X_0^+}(qaK_{X_0^+})\bigr)
\]
a finite module over that subalgebra. These finitely many residue
modules prove finite generation of the whole algebra.

Equation~\eqref{eq:family-canonical-algebra} expresses
$\mathcal B$ as the proper direct-image algebra of the powers
of the line bundle $L$. Applying \ref{an:finite-generation-lift}
together with \eqref{eq:family-algebra-base-change} proves
finite generation of $\mathcal B$ on a neighborhood of the
compact set $Z_0$. Properness of $\mathfrak Z\to D$ permits
another shrinking of $D$, after which we may define over the
entire restricted target
\begin{equation}\label{eq:total-flip-Proj}
 \mathfrak X^+=\Proj_{\mathfrak Z}\mathcal B.
\end{equation}

\subsection{Flatness and nearby canonical fibers}

By \eqref{eq:family-algebra-base-change}, the central fiber of
relative Proj is
\[
 (\mathfrak X^+)_0
 =\Proj_{Z_0}(\mathcal B/t\mathcal B)=X_0^+.
\]
Fix $c\in D$ and write $u=t-c$ for a local parameter at $c$.
Since $\mathfrak Z\to D$ is flat, multiplication by $u$ is
injective on the local rings of $\mathfrak Z$ along $Z_c$.
Each $\mathcal B_m$ is torsion-free on the normal space
$\mathfrak Z$, so multiplication by $u$ is also injective on
every graded piece. The same holds after homogeneous localization
and passage to the degree-zero part, and hence on the standard
local algebras of relative Proj. The analytic local rings of these
charts are flat over the corresponding local algebras of finite
type, so multiplication by $u$ remains injective after
analytification. They are therefore torsion-free, and hence flat,
over the discrete valuation ring $\OO_{D,c}$.
As $c$ was arbitrary, $\mathfrak X^+\to D$ is flat.
It is proper because $\mathfrak X^+\to\mathfrak Z$ is projective
and $\mathfrak Z\to D$ is proper.

The central fiber $X_0^+$ is the prescribed canonical space.
By \ref{an:canonical-deform}, after shrinking the disk the total
space $\mathfrak X^+$ is normal, canonical, and
$\Q$-Gorenstein, and all nearby fibers are canonical.
On the open set in \eqref{eq:family-common-open}, the canonical
sheaf is a line bundle, and \eqref{eq:total-flip-Proj} is the open
set itself. This open set meets every sufficiently nearby fiber.
Connectedness and irreducibility follow as before, so
$X_t^+\to Z_t$ is bimeromorphic for each nearby $t$.
Consequently $X_t\dashrightarrow X_t^+$ is bimeromorphic as well.

Since the central morphism $X_0^+\to Z_0$ is small, the same
argument with horizontal and vertical prime divisors proves that
$\mathfrak X^+\to\mathfrak Z$ is small. Choose a sufficiently
divisible integer $a$ such that the Proj sheaf $\OO(a)$ is a
relatively ample line bundle and
$aK_{\mathfrak X^+/D}$ is Cartier. On the common open set
containing all codimension-one points, $\OO(a)$ agrees with
$\omega_{\mathfrak X^+/D}^{[a]}(aX_0^+)$. Normality extends
this isomorphism to the whole space. The central fiber $X_0^+$
is the principal Cartier divisor defined by the parameter, so
twisting by it does not affect relative ampleness. Hence
$K_{\mathfrak X^+/D}$ is ample relative to
$\mathfrak X^+\to\mathfrak Z$. This completes the analytic
extension of the prescribed flip.

\subsection{Iteration of the finite sequence}

At each step, the construction preserves the proper flat family
and the prescribed canonical central fiber required for the next
step. We may therefore extend all the steps of
\eqref{eq:prescribed-central-program} successively. Since their
number $r$ is finite, the intersection of the finitely many
successive disk neighborhoods still contains a disk about the
origin. The central fiber of the final family is exactly
$X_0^{(r)}$. Its total space and nearby fibers have the asserted
singularity properties, and every nearby fiber is bimeromorphic
to the original fiber through the finite sequence of diagrams.
Lemma~\ref{lem:canonical-birational} gives
\eqref{eq:family-birational-plurigenera}, completing the proof of
Proposition~\ref{prop:finite-family}.

In particular, applying this proposition to the finite sequence
of Proposition~\ref{prop:finite-original} produces, from the
original family, a proper flat family
\begin{equation}\label{eq:terminal-family-output}
 f_{\min}:\mathfrak X_{\min}\longrightarrow D.
\end{equation}
Its central fiber is compact Kähler and terminal, with semiample
canonical divisor. Terminality implies
that the central fiber is smooth in codimension two, and a
projective resolution of a compact Kähler space is again compact
Kähler. The total space is canonical and $\Q$-Gorenstein, and
the nearby fibers are canonical. The resulting family therefore
satisfies the hypotheses of Theorem~\ref{thm:terminal-extension}.

\section{Extension of pluricanonical sections}\label{sec:terminal-extension}

For a smooth central fiber in Fujiki's class $\mathcal C$ and a
pluricanonical section with smooth reduced zero divisor, Levine's
cyclic-cover method expresses the extension obstruction as the trace
of an exterior derivative, whose
cohomology class vanishes by Hodge theory
\cite[Proposition~1.6]{Levine1983}. The same cover recovers the lower
pluricanonical powers \cite[Corollary~1.8]{Levine1983}.
His semiample-central-fiber result for smooth families with fibers in
Fujiki's class $\mathcal C$ is the smooth predecessor of the theorem
below \cite[Corollary~1.10]{Levine1983}.

For projective central fibers, Kollár--Mori proved the corresponding
invariance theorem for log-terminal singularities smooth in
codimension two and a globally generated locally free pluricanonical power
\cite[Theorem~12.5.5]{KM1992}.
This includes the projective case of the canonical statement below.
For terminal central fibers, Kollár records that the proof also applies
in Fujiki's class $\mathcal C$
\cite[comments following Conjecture~17]{Kollar2021}.

Cao and Păun further develop the infinitesimal method, proving a
lifting criterion under a weighted $L^2$ condition with a Kähler
central fiber \cite[Theorem~1.1]{CaoPaun2023}.

Theorem~\ref{thm:terminal-extension} establishes this class-$\mathcal C$
extension statement for canonical central fibers smooth in codimension
two, with semiample canonical divisor and a compact Kähler modification.
Its terminal special case suffices for the proof
of Theorem~\ref{thm:main}; the canonical formulation gives the
additional scope of the extension theorem.

We give an explicit logarithmic period proof based on the cyclic-cover
method of Levine and Kollár--Mori. Cohen--Macaulayness of the reflexive
canonical powers identifies the relevant first cohomology groups with
those of the smooth locus. Logarithmic Hodge theory on a compact Kähler
modification of the central cover, combined with a Leray comparison,
gives the intrinsic period injection used by the infinitesimal argument.

\begin{theorem}\label{thm:terminal-extension}
Let $f:\mathfrak X\to\Delta$ be a proper flat holomorphic map with connected
central fiber $X=X_0$ of dimension $\dim_\C X=n\ge3$. Suppose that, on a
neighborhood of $X$, the total space $\mathfrak X$ is normal,
$\Q$-Gorenstein, and has canonical singularities. Assume moreover that:
\begin{enumerate}[label=\textup{(\roman*)},leftmargin=*]
\item $X$ is normal, has canonical singularities, and satisfies
      $\operatorname{codim}_X\Sing X\ge3$;
\item there exists a proper bimeromorphic morphism
      $b:\widetilde X\to X$ from a compact Kähler manifold;
\item the canonical $\Q$-line bundle $K_X$ is semiample.
\end{enumerate}
Then, for every fixed integer $m\ge1$, every section of
$H^0(X,\omega_X^{[m]})$ extends, after shrinking the disk, to a holomorphic
section of
\[
 \cF_m:=\OO_{\mathfrak X}(mK_{\mathfrak X/\Delta}).
\]
After shrinking so that the nearby fibers are normal and have canonical
singularities, the function
\[
 t\longmapsto h^0(X_t,\omega_{X_t}^{[m]})
\]
is constant in a neighborhood of $0$.
\end{theorem}

The analytic deformation results in \ref{an:canonical-deform} guarantee
the asserted shrinking for the nearby fibers, together with canonical
singularities of the total space and the relative $\Q$-Gorenstein property
over a smooth disk.
A terminal central fiber satisfies \textup{(i)}. A compact Kähler central
fiber admits a projective resolution that is again compact Kähler, and hence
satisfies \textup{(ii)}. The conclusion is local for each fixed degree $m$.

\begin{proof}
Fix the pluricanonical degree $m$.
We choose a globally generated canonical power of degree $k\ge m$
and a general central section $s_0$. The character decomposition of
its cyclic cover will recover every canonical degree from $1$ to $k$.
Given a lift $s_r$ of $s_0$ to order $r$, we construct the corresponding
Artinian cover and prove a period injection for it and its truncations.
Levine's trace formula then annihilates the obstruction to lifting
$s_0$ to order $r+1$. This yields lifts to every finite order;
the character summands supply the required degree-$m$ lifts, and
coherence gives analytic extension.

\subsection{Canonical powers and cohomology on the smooth locus}

We first establish a local fact. If $Z$ is $\Q$-Gorenstein with canonical
singularities, then $\OO_Z(jK_Z)$ is maximal Cohen--Macaulay for
\emph{every integer} $j$, including negative integers.
Let $a$ be a local canonical index, and trivialize $aK_Z$ on a sufficiently
small neighborhood. The algebra of the corresponding index-one cover is
\[
 \bigoplus_{b=0}^{a-1}\OO_Z(-bK_Z).
\]
This finite cover is Gorenstein and has canonical singularities. More
explicitly, it is étale in codimension one, and the discrepancies of a
divisorial valuation and a lift of that valuation satisfy
\[
 a(E,Z')+1=e\bigl(a(F,Z)+1\bigr),\qquad e\ge1.
\]
Thus canonical singularities are preserved under this cover.
Canonical singularities are rational and therefore Cohen--Macaulay.
A system of parameters downstairs is a system of parameters in each of the
finite local rings upstairs, and hence is a regular sequence on the finite
direct image. Each character summand displayed above is a direct summand of
that module and is therefore Cohen--Macaulay. Tensoring with powers of the
invertible sheaf associated with $aK_Z$ proves the assertion for every
integer $j$.

Apply this observation to $\mathfrak X$, using the notation $\cF_j$ for
arbitrary integers $j$. The coordinate $t$ of the disk is a nonzerodivisor
on each torsion-free rank-one sheaf $\cF_j$, so $\cF_j$ is flat over the disk.
Set
\[
 A_r=\C[t]/(t^{r+1}),\qquad
 X_r=\mathfrak X\times_\Delta\Spec A_r,\qquad
 F_{j,r}=\cF_j/t^{r+1}\cF_j.
\]
Then $F_{j,r}$ is a Cohen--Macaulay module of dimension $n$ and is flat over
$A_r$.

Write $\Sigma=\Sing X$ and $U_r=X_r\setminus\Sigma$.
At every point of $U_0$, the flat morphism has a smooth fiber, so $f$ is
smooth on a neighborhood of that point. Consequently,
\begin{equation}\label{eq:ext-regular-powers}
 F_{j,r}|_{U_r}=\omega_{U_r/A_r}^{\otimes j}.
\end{equation}
In particular, $F_{j,0}$ is a rank-one Cohen--Macaulay module on the normal
space $X$. All of its associated primes are minimal; it is therefore
torsion-free, and it satisfies $S_2$. Thus it is
reflexive. The identification on the smooth locus in
\eqref{eq:ext-regular-powers} gives
\begin{equation}\label{eq:ext-reflexive-restriction}
 F_{j,0}=\omega_X^{[j]}.
\end{equation}
For every nearby normal canonical fiber $X_c$, multiplication by $t-c$
is again injective on $\cF_j$. Its quotient is a rank-one
Cohen--Macaulay sheaf on $X_c$, hence reflexive, and agrees with the
canonical power on the regular locus. Consequently,
\[
 \cF_j|_{X_c}=\omega_{X_c}^{[j]}.
\]
The ideal defining $\Sigma$ has grade at least three on every $F_{j,r}$.
Indeed, for a Cohen--Macaulay module, this grade is the infimum of the
localized depths at primes containing the ideal. These depths equal the
corresponding codimensions in the support, all of which are at least three.
The Artinian thickening does not change these codimensions. Local cohomology
therefore gives
\[
 \mathcal H^i_\Sigma(F_{j,r})=0\qquad(i<3).
\]
The local-to-global spectral sequence for cohomology with supports,
together with the long exact sequence for the complement, yields
\begin{equation}\label{eq:ext-depth-base}
 H^1(X_r,F_{j,r})
 \xrightarrow{\ \sim\ }
 H^1(U_r,\omega_{U_r/A_r}^{\otimes j}).
\end{equation}
The same grade computation gives the Hartogs extension property in
codimension two used below.

Since $\mathfrak X$ is $\Q$-Gorenstein and $X$ is compact, there is a common
Cartier index on a neighborhood of $X$. Taking a common multiple with a
semiample exponent on the central fiber, choose
$k\ge\max\{m,2\}$ such that $\cF_k$ is invertible and
$F_{k,0}=\OO_X(kK_X)$ is generated by its global sections.

\subsection{The central canonical cyclic cover}

We choose a section $s_0\in H^0(X,kK_X)$ whose canonical root cover
is normal, Gorenstein, and canonical, and whose inverse image of
$X_{\mathrm{reg}}$ is smooth.
We follow the canonical-cover construction of Kawamata recalled in
\cite[Lemma~12.5.10]{KM1992}, giving the Gorenstein and analytic
Bertini details used below.

Consider the analytic canonical line space
\begin{equation}\label{eq:ext-canonical-line-space}
 p:Z=\Spec_X\left(\bigoplus_{j\ge0}\OO_X(-jK_X)\right)\longrightarrow X.
\end{equation}
This algebra is locally of finite type: after trivializing $aK_X$, its
degree-$j+a$ part is obtained from its degree-$j$ part by multiplication
with a local generator in degree $a$. Over $X_{\mathrm{reg}}$, the space
$Z$ is the usual total space of $K_X$.

Locally, take the index-one cover $\widehat X\to X$ with group $\mu_a$,
and choose a canonical generator $\theta$ with character $\chi$. Then
\begin{equation}\label{eq:ext-line-quotient}
 Z\simeq(\widehat X\times\C_z)/\mu_a,\qquad
 g(x,z)=\bigl(gx,\chi(g)^{-1}z\bigr).
\end{equation}
Indeed, the coefficient of $z^j$ in an invariant polynomial has character
$\chi^j$, which gives precisely $\OO_X(-jK_X)$.
The top form $\theta\wedge dz$ is invariant. The finite quotient is normal
and Cohen--Macaulay, and this invariant generator trivializes the dualizing
sheaf on the codimension-one locus. Reflexive extension identifies that
sheaf with an invertible trivial sheaf, so $Z$ is Gorenstein.
The quotient is étale in codimension one, and
$\widehat X\times\C$ has canonical singularities.
The discrepancy formula implies that $Z$ is klt. Since $K_Z$ is Cartier,
its exceptional discrepancies are integers, hence nonnegative; thus $Z$
is canonical.

Let $W=H^0(X,kK_X)$, and denote by $\tau_k$ the degree-$k$ tautological
section of $p^*(kK_X)$. The finite-dimensional system
\begin{equation}\label{eq:ext-bpf-system}
 \C\tau_k+p^*W
\end{equation}
is basepoint-free on all of $Z$: at a point lying over $x$, some section
of $W$ is nonzero at $x$. Take a log resolution of $Z$ and write
\[
 K_{\widetilde Z}=g^*K_Z+\sum_i a_iE_i,\qquad a_i\ge0.
\]
The pullback of \eqref{eq:ext-bpf-system} remains basepoint-free.
A general member is smooth on $\widetilde Z$, is transverse to every
exceptional stratum, and contains no exceptional center.
These assertions follow by applying Sard's theorem to the projection
from the smooth incidence space and then applying the same argument on
each stratum. Even if $\widetilde Z$ is noncompact, there are at most
countably many strata; the bad parameter values form a countable union
of measure-zero sets and can therefore be avoided simultaneously.
We may also avoid the hyperplane on which the coefficient of $\tau_k$
vanishes, and arrange that the zero divisor of the resulting $s_0$ is
smooth and reduced on $X_{\mathrm{reg}}$.

After normalizing the coefficient of $\tau_k$, the member has the form
\begin{equation}\label{eq:ext-central-root}
 Y_0=\{\tau_k-p^*s_0=0\}.
\end{equation}
Its total transform on the resolution has no exceptional divisorial
component. As a Cartier hypersurface in the Cohen--Macaulay space $Z$,
the space $Y_0$ is Cohen--Macaulay. Bertini's theorem, together with the
requirement that no component of $\Sing Z$ be contained in $Y_0$, shows
that it is smooth in codimension one. Hence $Y_0$ is normal.
Adjunction on its smooth strict transform gives the nonnegative
exceptional coefficients $a_i$ above, so $Y_0$ is canonical.
It is also Gorenstein, being a Cartier hypersurface in a Gorenstein space.

The monic relation of degree $k$ defines the finite cover
\begin{equation}\label{eq:ext-central-algebra}
 q_0:Y_0\to X,\qquad
 (q_0)_*\OO_{Y_0}
 =\bigoplus_{a=0}^{k-1}\OO_X(-aK_X).
\end{equation}
Over $X_{\mathrm{reg}}$, its equation is $w^k=s_0$.
Since the branch divisor is smooth and reduced,
$q_0^{-1}(X_{\mathrm{reg}})$ is smooth, including at the ramification
points. Set
\[
 Z_0=q_0^{-1}\Sigma,\qquad V_0=Y_0\setminus Z_0.
\]
Finiteness gives $\operatorname{codim}_{Y_0}Z_0\ge3$.

\subsection{Cyclic covers over Artinian thickenings}

Suppose that a section
$s_r\in H^0(X_r,F_{k,r})$ has already been constructed, with central
restriction equal to the section $s_0$ chosen above. Define
\begin{equation}\label{eq:ext-artin-algebra}
 q_r:Y_r=\Spec_{X_r}
 \left(\bigoplus_{a=0}^{k-1}F_{-a,r}\right)\longrightarrow X_r.
\end{equation}
Multiplication in degrees reaching $k$ is reduced using $s_r$.
The multiplication maps are induced by the canonical-power multiplication
on the total space and the invertibility of $\cF_k$, and remain defined
after reduction modulo $t^{r+1}$. On $U_r$, this is the usual algebra
with equation $w^k=s_r$. Associativity may also be checked on that open
set and then extended by the Hartogs property of the target modules.
Thus \eqref{eq:ext-artin-algebra} is a well-defined global algebra.

Each summand is flat over $A_r$ and maximal Cohen--Macaulay. These
properties of the finite direct image imply that $Y_r$ is flat over
$A_r$ and Cohen--Macaulay.
Let $V_r=Y_r\setminus Z_0$, identifying the underlying topological spaces
of the thickenings with their central fibers.
Then $V_r/A_r$ is smooth. Away from the branch locus, $kw^{k-1}$ is a
unit. At a branch point, a fiber-coordinate derivative of $s_0$ is a
unit and remains a unit on the thickening. Replacing that coordinate by
$s_r$ puts the equation into the form $z_1=w^k$.

Set
\begin{equation}\label{eq:ext-top-line}
 L_r=q_r^*F_{k,r}.
\end{equation}
This is an invertible sheaf. The total space of the line bundle
$\omega_{U_r/A_r}$ has trivial canonical bundle relative to $A_r$:
the canonical factor from the base cancels the inverse factor from the
line-bundle direction. Adjunction for the degree-$k$ equation gives
\begin{equation}\label{eq:ext-adjunction}
 L_r|_{V_r}=\omega_{V_r/A_r}.
\end{equation}
All these identifications commute with truncation.
The ideal of the removed set has grade at least three on the invertible
sheaf $L_r$ on the Cohen--Macaulay space $Y_r$. Hence
\begin{equation}\label{eq:ext-depth-cover}
 H^1(Y_r,L_r)\xrightarrow{\ \sim\ }
 H^1(V_r,\omega_{V_r/A_r}).
\end{equation}
Since $F_{k,r}$ is invertible, the projection formula gives
\begin{equation}\label{eq:ext-top-direct-sum}
 (q_r)_*L_r
 =\bigoplus_{a=0}^{k-1}F_{k-a,r}
 =\bigoplus_{\ell=1}^{k}F_{\ell,r}.
\end{equation}
For each existing $s_r$, \eqref{eq:ext-artin-algebra} thus defines a cover
over $A_r$, compatibly with its truncations.

\subsection{The logarithmic period injection}

We first construct a compact Kähler resolution of $Y_0$.
Choose the compact Kähler resolution $b:\widetilde X\to X$
in hypothesis \textup{(ii)}.
Since the canonical discrepancies are nonnegative, the pullback of
$s_0$ can be regarded as a holomorphic section of $kK_{\widetilde X}$.
Its root hypersurface is a compact subspace of
$\mathbf P(\OO_{\widetilde X}\oplus K_{\widetilde X})$;
the monic equation ensures that it does not meet the divisor at infinity.
Its algebra is locally free over $\widetilde X$ and generically separable,
so the hypersurface is reduced. A projective embedded resolution is
compact Kähler. Over the dense open set where $b$ is an isomorphism,
this root cover agrees with $Y_0$.

The resulting map to the normal finite cover $Y_0$ extends over the entire
resolution. Indeed, locally choose integral generators for the finite
algebra defining $Y_0$ over the base. On the dense open set where the map
is defined, their pullbacks satisfy monic equations with locally bounded
coefficients, and the roots are consequently locally bounded.
The removable-singularity theorem on a normal space extends them
holomorphically. Their defining equations and gluing identities continue
from the dense open set. The resulting map is proper and bimeromorphic
on each component. Further projective blowups principalize the
inverse-image ideal of $Z_0$ and give
\[
 \widetilde q:\widetilde Y\to Y_0,\qquad
 E=(\widetilde q^{-1}Z_0)_{\mathrm{red}},
\]
where $\widetilde Y$ is compact Kähler and $E$ is a simple normal crossing
divisor. Set
\[
 W=\widetilde Y\setminus E,\qquad
 h=\widetilde q|_W:W\longrightarrow V_0.
\]
The morphism $h$ is a proper modification between smooth spaces.

Since $Y_0$ is normal and Gorenstein,
\eqref{eq:ext-adjunction} extends reflexively to $L_0=\omega_{Y_0}$.
Canonical singularities give inclusions
\[
 \widetilde q^*L_0
 \hookrightarrow \omega_{\widetilde Y}
 \hookrightarrow \Omega^n_{\widetilde Y}(\log E).
\]
We claim that the induced map
\[
 H^1(Y_0,L_0)\longrightarrow
 H^1\bigl(\widetilde Y,\Omega^n_{\widetilde Y}(\log E)\bigr)
\]
is injective. The sheaf identity in
Lemma~\ref{lem:canonical-birational} gives
$h_*\omega_W=\omega_{V_0}$. The low-degree Leray sequence therefore
gives an injection
\[
 H^1(V_0,\omega_{V_0})\hookrightarrow H^1(W,\omega_W).
\]
It is induced by pullback of top differentials. Restricting the claimed
map to $W$ gives the composite of this injection with the isomorphism
\eqref{eq:ext-depth-cover} for $r=0$. This proves the claim.

For the compact Kähler simple normal crossing pair $(\widetilde Y,E)$,
we use logarithmic Hodge-to-de Rham degeneration in the compact Kähler
formulation of Liu, Rao, and Wan \cite[Theorem~3.2]{LRW2019}.
The spectral sequence degenerates at $E_1$.
Since there are no logarithmic forms of degree greater than $n$, the
edge map from the top column,
\[
 H^1\bigl(\widetilde Y,\Omega^n_{\widetilde Y}(\log E)\bigr)
 \longrightarrow
 \mathbb H^{n+1}\bigl(\widetilde Y,\Omega^\bullet_{\widetilde Y}(\log E)\bigr),
\]
is injective. Logarithmic de Rham comparison
\cite[Section~3.2]{LRW2019} identifies the target with
$H^{n+1}(W,\C)$. We have therefore obtained an injection from
$H^1(Y_0,L_0)$ to $H^{n+1}(W,\C)$.

Define $\per_0$ by restriction to $V_0$, inclusion of its top
holomorphic de Rham term, and holomorphic de Rham comparison.
Naturality under restriction and pullback identifies the preceding
injective composite with
\[
 H^1(Y_0,L_0)\xrightarrow{\ \per_0\ }
 H^{n+1}(V_0,\C)\xrightarrow{\ h^*\ }H^{n+1}(W,\C).
\]
Consequently, the central period map is injective:
\begin{equation}\label{eq:ext-central-period}
 \per_0:H^1(Y_0,L_0)\hookrightarrow H^{n+1}(V_0,\C).
\end{equation}

\subsection{Period injections over Artinian thickenings}

Fix an existing lift $s_r$ and its compatible truncations.
For $0\le j\le r$, the morphism $V_j/A_j$ is smooth.
The relative holomorphic Poincaré lemma gives a canonical
quasi-isomorphism
\[
 \underline{A_j}\simeq \Omega^\bullet_{V_j/A_j}.
\]
In local smooth coordinates this is the ordinary Poincaré lemma applied
coefficient by coefficient in $t$. The canonical inclusions of the constant
sheaf identify these local quasi-isomorphisms on overlaps. Consequently,
\[
 \mathbb H^{n+1}(V_j,\Omega^\bullet_{V_j/A_j})
 =H^{n+1}(V_0,\C)\otimes_\C A_j.
\]
Restriction to $V_j$, followed by inclusion of the top de Rham term,
defines
\begin{equation}\label{eq:ext-artin-period}
 \per_j:H^1(Y_j,L_j)\longrightarrow H^{n+1}(V_0,\C)\otimes_\C A_j.
\end{equation}
These maps commute both with reduction modulo $t$ and with multiplication
by $t$.

We prove by induction that $\per_j$ is injective and that
\begin{equation}\label{eq:ext-top-section-surjection}
 H^0(Y_j,L_j)\longrightarrow H^0(Y_0,L_0)
 \quad\text{is surjective}.
\end{equation}
The initial period injection is \eqref{eq:ext-central-period}, and the
assertion about sections is automatic for $j=0$.
For $j\ge1$, regard $L_{j-1}$ and $L_0$ as sheaves on $Y_j$ through
their natural closed embeddings. Flatness gives an exact sequence
\begin{equation}\label{eq:ext-t-exact}
 0\longrightarrow L_{j-1}\xrightarrow{\ t\ }L_j
 \longrightarrow L_0\longrightarrow0.
\end{equation}
Assume that $\per_{j-1}$ is injective.
If $b\in H^1(Y_{j-1},L_{j-1})$ maps to zero under multiplication by $t$,
compatibility gives $t\,\per_{j-1}(b)=0$.
Multiplication $t:A_{j-1}\to A_j$ is injective, so
$\per_{j-1}(b)=0$, and hence $b=0$.
Thus the multiplication-by-$t$ map on $H^1$ in
\eqref{eq:ext-t-exact} is injective. The long exact sequence gives
\eqref{eq:ext-top-section-surjection}.

Now let $a\in H^1(Y_j,L_j)$ have zero period.
Its reduction to the central fiber vanishes by
\eqref{eq:ext-central-period}. Hence $a=tb$ for some
$b\in H^1(Y_{j-1},L_{j-1})$.
The identity $0=\per_j(a)=t\,\per_{j-1}(b)$ and the same injectivity
imply $b=0$. This proves the injectivity of $\per_j$ and completes the
induction.

In particular,
\begin{equation}\label{eq:ext-d-zero}
 \dd:H^1(V_j,\Omega^{n-1}_{V_j/A_j})
 \longrightarrow H^1(V_j,\omega_{V_j/A_j})
 \quad\text{is zero}.
\end{equation}
Indeed, applying $\dd$ to a \v{C}ech cocycle gives a boundary in the
total \v{C}ech--de Rham complex, so its de Rham period is zero.
By \eqref{eq:ext-depth-cover}, this image comes from a unique class in
$H^1(Y_j,L_j)$. The injectivity of \eqref{eq:ext-artin-period} then makes
it zero.

\subsection{Levine's obstruction and trace formulas}

We apply Levine's local obstruction calculation
\cite[Lemmas~1.2--1.5, pp.~295--298]{Levine1983}.
We give the formulas in the present notation, including the identification
of the central-section obstruction and the holomorphic map across the
ramification divisor.

Suppose again that $s_r$ has been constructed. Work on the smooth open
part of the original family, with fiber coordinates $(z_i,t)$ and
transition functions
\[
 z_i=g_{ij}(z_j,t),\qquad
 J_{ij}=\det\left(\frac{\partial g_{ij}}{\partial z_j}\right),\qquad
 \rho_{ij}=\sum_\alpha
 \frac{\partial g_{ij}^{\alpha}}{\partial t}
 \frac{\partial}{\partial z_i^\alpha}.
\]
Write the section as $s_i(dz_i)^k$. Its transition rule is
\[
 s_i(g_{ij},t)=J_{ij}^{-k}s_j\pmod {t^{r+1}}.
\]
Differentiating the transition rule gives the cocycle, expressed in
the $i$-coordinates,
\begin{equation}\label{eq:ext-Levine-D}
 D(s_r)_{ij}
 =\left(k\,\operatorname{div}(\rho_{ij})s_i+\rho_{ij}(s_i)\right)(dz_i)^k.
\end{equation}
Let $\delta_{r+1}(s_0)$ be the connecting class associated with
\begin{equation}\label{eq:ext-central-obstruction-sequence}
 0\longrightarrow F_{k,r}\xrightarrow{\ t\ }F_{k,r+1}
 \longrightarrow F_{k,0}\longrightarrow0.
\end{equation}
Then
\begin{equation}\label{eq:ext-Levine-obstruction}
 [D(s_r)]=(r+1)\delta_{r+1}(s_0)|_{U_r}.
\end{equation}
Choose arbitrary next-order lifts of the local coefficients of $s_r$
and adopt the \v{C}ech convention
\[
 J_{ij}^{\,k}s_i(g_{ij},t)-s_j=t^{r+1}\zeta_{ij}
 \pmod {t^{r+2}}
\]
for the \v{C}ech difference.
Dividing this difference by $t$ gives the representative
$t^r\zeta_{ij}$ of the connecting class in
\eqref{eq:ext-central-obstruction-sequence}.
Differentiation by $\partial_t$ gives $(r+1)t^r\zeta_{ij}$ on the right.
On the left, the differences of the local $\partial_t s_i$ form a
coboundary, and the remaining term is
\eqref{eq:ext-Levine-D}, using
$\partial_t J_{ij}/J_{ij}=\operatorname{div}(\rho_{ij})$.
This proves \eqref{eq:ext-Levine-obstruction}.
The class in \eqref{eq:ext-Levine-obstruction} is the connecting image of
the central section $s_0$; the chosen finite-order lift $s_r$ is used to
compute it.

On $V_r$, write the tautological pulled-back canonical section as
$\bar\omega=w\,dz$, and set
\[
 \mu_r=\langle\rho,\bar\omega\rangle
 \in H^1(V_r,\Omega^{n-1}_{V_r/A_r}).
\]
Differentiating the transition identity
$g_{i\ell}=g_{ij}\circ g_{j\ell}$ with respect to $t$ gives the
\v{C}ech cocycle identity for $\rho$.
The root coordinates satisfy $w_i\,dz_i=w_j\,dz_j$, so
$\bar\omega$ is a global section. The contraction therefore defines
a cohomology class of $(n-1)$-forms.
The contraction is first taken in the pullbacks from $U_r$ of the
relative tangent sheaf and canonical line, and is then composed with
the natural differential map to forms on $V_r$.
The local equation gives
\[
 kw^{k-1}\dd w=\dd s_i.
\]
Off the ramification locus, the differential identifies top forms
with the pullback of the canonical line of $U_r$. Under this
identification, multiplication by $\bar\omega^{k-1}$ defines a map
to $q_r^*(\omega_{U_r/A_r}^{\otimes k})$.
This map extends holomorphically across the ramification locus,
as the following local calculation shows.
Choose a coordinate index for which $\partial s_i/\partial z_1$
is a unit. A local frame of the canonical line of the cover is
\begin{equation}\label{eq:ext-branch-frame}
 \frac{dz}{kw^{k-1}}
 =\left(\frac{\partial s_i}{\partial z_1}\right)^{-1}
 \dd w\wedge dz_2\wedge\cdots\wedge dz_n.
\end{equation}
Multiplication by $(w\,dz)^{k-1}$ sends this frame to $(dz)^k/k$.
Thus the operation is a holomorphic sheaf map, including at the
ramification locus.

Expanding
$\dd(w\,\iota_\rho dz)=\dd w\wedge\iota_\rho dz
+w\,\operatorname{div}(\rho)\,dz$
and using the root equation gives
\begin{equation}\label{eq:ext-trace-local}
 \bar\omega^{k-1}\dd\mu_r
 =\left(\operatorname{div}(\rho)s_i+
          \frac1k\rho(s_i)\right)(dz_i)^k.
\end{equation}
This expression is invariant under the cyclic group.
The monic finite cover has total degree $k$, so its trace is $k$ times
this invariant expression. Consequently,
\begin{equation}\label{eq:ext-trace-obstruction}
 [D(s_r)]
 =\Tr_{V_r/U_r}\left(\bar\omega^{k-1}\dd\mu_r\right).
\end{equation}
The same computation applies if the cover is disconnected: the monic
algebra still has total rank $k$.
These calculations use only smooth coordinate charts and apply on the
smooth open set considered here in every dimension.

By \eqref{eq:ext-d-zero}, the cohomology class of $\dd\mu_r$ vanishes.
The holomorphic sheaf map and the trace then make
\eqref{eq:ext-trace-obstruction} zero.
Combining \eqref{eq:ext-Levine-obstruction} with
\eqref{eq:ext-depth-base}, and using that $r+1$ is invertible over $\C$,
gives
$\delta_{r+1}(s_0)=0$ in $H^1(X_r,F_{k,r})$.
Thus $s_0$ has at least one lift to order $r+1$.
Starting with $r=0$, induction therefore shows that, for every $r\ge0$,
there exists $s_r\in H^0(X_r,F_{k,r})$ whose restriction to $X$ is $s_0$.

\subsection{Recovery of the lower pluricanonical powers}

The recovery of lower powers is the cyclic-cover argument of
\cite[Corollary~1.8]{Levine1983}, whose proof Levine credits to
Shepherd-Barron. Here we apply the character decomposition over each
fixed Artinian base.

Fix a target order $r$, and choose one lift $s_r$ constructed above
together with its own truncations.
Apply \eqref{eq:ext-top-section-surjection} to this finite chain and
use \eqref{eq:ext-top-direct-sum}. It follows that
\[
 \bigoplus_{\ell=1}^{k}H^0(X_r,F_{\ell,r})
 \longrightarrow
 \bigoplus_{\ell=1}^{k}H^0(X,F_{\ell,0})
 \quad\text{is surjective}.
\]
The restriction map respects the direct-sum decomposition.
Since $m\le k$, every section of $H^0(X,F_{m,0})$ therefore lifts
to order $r$.
Since the target order $r$ was arbitrary, this proves surjectivity at
every finite order.

\subsection{Passage from finite-order lifts to analytic sections}

The passage from lifting to every finite order to analytic local
constancy is already present, for line bundles in smooth families,
in \cite[Lemma~1.1]{Levine1983}. For coherent analytic sheaves flat
over the base of a proper flat family over a smooth Stein curve,
the relation between extension, base change, and torsion freeness is
developed in \cite[Theorem~1.2 and Appendix~A]{FLR2026}.
We give a direct finite-module argument for the reflexive canonical
powers, using the finite-order lifts constructed above.

Set
\[
 R=\C\{t\},\qquad E=(f_*\cF_m)_0,\qquad
 G=(R^1f_*\cF_m)_0.
\]
Proper coherence makes $E$ and $G$ finite $R$-modules.
The bounded exponent of the torsion in $G$ will make one sufficiently
high finite-order lift enough for analytic extension.
Apply direct image to the exact sequence given by multiplication
by $t^a$ on $\cF_m$, and take the stalk at $0$. The result is
\begin{equation}\label{eq:ext-DVR-exact}
 0\longrightarrow E/t^aE
 \longrightarrow H^0(X_{a-1},F_{m,a-1})
 \xrightarrow{\ \partial_a\ }G[t^a]\longrightarrow0,
\end{equation}
where $G[t^a]=\ker(t^a:G\to G)$.
In the map from the sequence modulo $t^a$ to the sequence modulo $t$,
the middle map on $\cF_m$ is the identity and the left map on
$\cF_m$ is multiplication by $t^{a-1}$.
Naturality gives
\begin{equation}\label{eq:ext-DVR-naturality}
 \partial_1\circ\operatorname{res}=t^{a-1}\partial_a.
\end{equation}

The torsion submodule of the finite $R$-module $G$ is annihilated
by some fixed power $t^B$.
Given a central section, choose $a>B$ and use the lift to order
$a-1$ supplied by the preceding step.
Its image under $\partial_a$ lies in $G[t^a]$, hence in the torsion
submodule. Equation \eqref{eq:ext-DVR-naturality} therefore makes
the $\partial_1$-image of the central section zero.
By \eqref{eq:ext-DVR-exact}, the section comes from $E/tE$,
and consequently from a holomorphic section germ on the original
analytic family.
Choosing a basis of the finite-dimensional central
section space gives a common smaller disk on which all its elements
extend.

Since $\cF_m$ is flat over the disk, its direct image $f_*\cF_m$
is torsion-free and hence locally free.
The surjectivity just proved identifies its central fiber with
the entire space $H^0(X,F_{m,0})$.
At every nearby point, the analogous exact sequence injects the fiber
of this locally free direct image into the corresponding fiber
cohomology space $H^0$. This gives the central dimension as a lower
bound. Proper upper semicontinuity gives the central dimension as an
upper bound for nearby fibers. The two bounds agree, proving local
constancy of $h^0$.
Finally, \eqref{eq:ext-reflexive-restriction} and its counterpart on
nearby fibers identify these dimensions with the required reflexive
plurigenera.
\end{proof}

\begin{remark}
The codimension-three condition gives the restriction isomorphisms
\eqref{eq:ext-depth-base} and \eqref{eq:ext-depth-cover}.
The compact Kähler resolution is used to establish the central
period injection; the Artinian period maps are then constructed on
the smooth open thickenings $V_j/A_j$.
Normality and canonical singularities identify the reflexive
pluricanonical powers on nearby fibers.
\end{remark}

\section{Proof of the main theorem}
\label{sec:main-proof}

\subsection{Fibers of dimension at most two}

If the fibers have dimension zero, connectedness makes each of them a
single point. Hence $P_m(X_t)=1$ for every $m\ge1$.

For one-dimensional fibers, Ehresmann's theorem makes the genus $g$
locally constant. Riemann--Roch and Serre duality give, for a smooth
compact connected curve $C$ of genus $g$,
\[
 P_m(C)=
 \begin{cases}
  0,&g=0,\\
  1,&g=1,\\
  g,&g\ge2\text{ and }m=1,\\
  (2m-1)(g-1),&g\ge2\text{ and }m\ge2.
 \end{cases}
\]
Indeed, $K_C$ has negative degree when $g=0$ and is trivial when $g=1$.
For $g\ge2$ and $m\ge2$, the line bundle $K_C^{\otimes(1-m)}$ has
negative degree, so its space of sections vanishes. The displayed
formula proves local constancy of every plurigenus.

For two-dimensional fibers, the assertion is Iitaka's deformation
invariance theorem for compact complex surfaces
\cite[Theorem~III, p.~247]{Iitaka1970}.
Since $\Delta$ is connected, these local constancy statements prove
Theorem~\ref{thm:main} in dimensions at most two.

\subsection{Fibers of dimension at least three}

Assume now that $n\ge3$. Fix an integer $m\ge1$ and a point
$t_0\in\Delta$. After translating the parameter and restricting to a
smaller disk, write $t_0=0$.

If $P_m(X_0)=0$, upper semicontinuity for the line bundle
$K_{\cX/\Delta}^{\otimes m}$ gives $P_m(X_t)=0$ for all sufficiently
small $t$.

Suppose that $P_m(X_0)>0$. A nonzero $m$-canonical section defines an
effective divisor $D_s$ with $mK_{X_0}\sim D_s$.
The positive current $[D_s]/m$ shows that $K_{X_0}$ is pseudoeffective.
Proposition~\ref{prop:finite-original} therefore supplies a finite
ordinary $K$-MMP
\begin{equation}\label{eq:final-central-program}
 X_0=X_0^{(0)}\dashrightarrow X_0^{(1)}\dashrightarrow\cdots
 \dashrightarrow X_0^{(N)}=:M_0
\end{equation}
whose models are compact Kähler and terminal, and whose endpoint has
semiample canonical divisor.

Proposition~\ref{prop:finite-family} extends
\eqref{eq:final-central-program} over a smaller disk $D\ni0$.
Write the resulting family as
\[
 f:\mathfrak M\longrightarrow D.
\]
It is proper and flat, with central fiber $M_0$.
After shrinking $D$, its total space near the central fiber is normal,
canonical, and $\Q$-Gorenstein, and each fiber $M_t$ is normal canonical
and bimeromorphic to $X_t$.

Terminal singularities are smooth in codimension two, so
\[
 \operatorname{codim}_{M_0}\Sing M_0\ge3.
\]
A projective resolution of the compact Kähler space $M_0$ is compact
Kähler. Theorem~\ref{thm:terminal-extension} therefore gives
\[
 h^0(M_t,\omega_{M_t}^{[m]})
 =h^0(M_0,\omega_{M_0}^{[m]})
\]
for sufficiently small $t$.
By Lemma~\ref{lem:canonical-birational},
\begin{align*}
 P_m(X_t)
 &=h^0(M_t,\omega_{M_t}^{[m]})\\
 &=h^0(M_0,\omega_{M_0}^{[m]})
 =P_m(X_0).
\end{align*}
Thus $P_m$ is locally constant on $\Delta$ for every $m\ge1$.
Since the disk is connected, Theorem~\ref{thm:main} follows.

\hfill$\square$

\section*{AI Disclosure}
All mathematical derivations and proofs were produced by AI. The human
author formulated the problem, provided relevant literature, and guided
the research direction.
\begin{thebibliography}{99}

\bibitem{Boucksom}
S.~Boucksom,
\emph{Divisorial Zariski decompositions on compact complex manifolds},
Ann. Sci. Éc. Norm. Supér. (4) \textbf{37} (2004), no.~1, 45--76.
\href{https://doi.org/10.1016/j.ansens.2003.04.002}{doi:10.1016/j.ansens.2003.04.002}.

\bibitem{CaoPaun2023}
J.~Cao and M.~Păun,
\emph{Infinitesimal extension of pluricanonical forms},
preprint, 9 December 2020, revised 28 March 2023.
\href{https://arxiv.org/abs/2012.05063v5}{arXiv:2012.05063v5}.

\bibitem{CRW2026}
J.~Chen, S.~Rao, and K.~Wang,
\emph{Matsumura's extension problem for pluricanonical forms in
Kähler families I: the smooth and essentially Moishezon cases},
preprint, 30 September 2026.
\href{https://arxiv.org/abs/2609.40040v1}{arXiv:2609.40040v1}.

\bibitem{DNWZ2026}
F.~Deng, J.~Ning, Z.~Wang, and X.~Zhou,
\emph{Invariance of plurigenera for Kähler families with nef canonical bundles},
preprint, 29 September 2026.
\href{https://arxiv.org/abs/2609.37804v1}{arXiv:2609.37804v1}.

\bibitem{FLR2026}
J.~Feng, Y.~Li, and S.~Rao,
\emph{Aspects of deformation invariance of plurigenera of Kähler varieties},
preprint, 13 October 2026.
\href{https://arxiv.org/abs/2610.17817v1}{arXiv:2610.17817v1}.

\bibitem{FujinoVanishing2025}
O.~Fujino,
\emph{Vanishing theorems for projective morphisms between complex analytic spaces},
Math. Res. Lett. \textbf{32} (2025), no.~3, 739--769.
\href{https://www.math.kyoto-u.ac.jp/~fujino/analytic-vanishing18.pdf}{Author's version}.

\bibitem{HLR2026}
C.~D.~Hacon, Y.~Li, and S.~Rao,
\emph{On Pseudo-Effectivity and Volumes of Adjoint Classes in Kähler
Families with Projective Central Fiber},
preprint, 4 February 2026, revised 5 March 2026.
\href{https://arxiv.org/abs/2602.04158v2}{arXiv:2602.04158v2}.

\bibitem{HX2026}
C.~Hacon and L.~Xie,
\emph{On the Kähler MMP and the transcendental base-point-free theorem},
preprint, 27 July 2026.
\href{https://arxiv.org/abs/2607.24986v1}{arXiv:2607.24986v1}.

\bibitem{Hironaka1975}
H.~Hironaka,
\emph{Flattening theorem in complex-analytic geometry},
Amer. J. Math. \textbf{97} (1975), no.~2, 503--547.
\href{https://doi.org/10.2307/2373721}{doi:10.2307/2373721}.

\bibitem{Iitaka1970}
S.~Iitaka,
\emph{Deformations of compact complex surfaces II},
J. Math. Soc. Japan \textbf{22} (1970), no.~2, 247--261.
\href{https://doi.org/10.2969/jmsj/02220247}{doi:10.2969/jmsj/02220247}.

\bibitem{Kollar2021}
J.~Kollár,
\emph{Deformations of varieties of general type},
Milan J. Math. \textbf{89} (2021), 345--354.
\href{https://arxiv.org/abs/2101.10986v3}{arXiv:2101.10986v3}.

\bibitem{KM1992}
J.~Kollár and S.~Mori,
\emph{Classification of three-dimensional flips},
J. Amer. Math. Soc. \textbf{5} (1992), no.~3, 533--703.
\href{https://doi.org/10.1090/S0894-0347-1992-1149195-9}{doi:10.1090/S0894-0347-1992-1149195-9}.

\bibitem{Levine1983}
M.~Levine,
\emph{Pluri-canonical divisors on Kähler manifolds},
Invent. Math. \textbf{74} (1983), 293--303.
\href{https://doi.org/10.1007/BF01394318}{doi:10.1007/BF01394318}.

\bibitem{LPWZ2026}
Y.~Li, M.~Păun, Z.~Wang, and X.~Zhou,
\emph{Psh variation of singular Kähler--Einstein metrics and
Invariance of plurigenera for smooth Kähler families},
preprint, 6 October 2026.
\href{https://arxiv.org/abs/2610.07826v1}{arXiv:2610.07826v1}.

\bibitem{LRW2019}
K.~Liu, S.~Rao, and X.~Wan,
\emph{Geometry of logarithmic forms and deformations of complex structures},
J. Algebraic Geom. \textbf{28} (2019), 773--815.
\href{https://doi.org/10.1090/jag/723}{doi:10.1090/jag/723}.

\bibitem{LW2026}
K.~Liu and X.~Wan,
\emph{Invariance of logarithmic plurigenera and adjoint volumes
in Kähler families},
preprint, 19 October 2026.
\href{https://arxiv.org/abs/2610.26519v1}{arXiv:2610.26519v1}.

\bibitem{Nakayama2004}
N.~Nakayama,
\emph{Zariski-decomposition and abundance},
MSJ Memoirs, vol.~14, Mathematical Society of Japan, Tokyo, 2004.

\bibitem{OAI033}
OpenAI,
\href{https://github.com/openai/math/tree/4eaee87548250b2682affda7d05d947c05554bae/preprints/Orbifold-and-logarithmic-Iitaka-subadditivity-September-26-2026}{\emph{Orbifold and logarithmic Iitaka subadditivity}},
OAI~033, preprint, 26 September 2026.

\bibitem{OAI034}
OpenAI,
\href{https://github.com/openai/math/tree/4eaee87548250b2682affda7d05d947c05554bae/preprints/Log-abundance-for-compact-Kahler-spaces-under-logarithmic-Iitaka-subadditivity-October-4-2026}{\emph{Log abundance for compact Kähler spaces under logarithmic Iitaka subadditivity}},
OAI~034, preprint, 4 October 2026.

\bibitem{Siu1998}
Y.-T.~Siu,
\emph{Invariance of plurigenera},
Invent. Math. \textbf{134} (1998), 661--673.
\href{https://doi.org/10.1007/s002220050276}{doi:10.1007/s002220050276}.

\bibitem{Siu2002}
Y.-T.~Siu,
\emph{Extension of twisted pluricanonical sections with plurisubharmonic
weight and invariance of semipositively twisted plurigenera for manifolds
not necessarily of general type},
in \emph{Complex Geometry} (Göttingen, 2000),
Springer, Berlin, 2002, pp.~223--277.
\href{https://doi.org/10.1007/978-3-642-56202-0_15}{doi:10.1007/978-3-642-56202-0\_15}.

\end{thebibliography}

\end{document}
